\documentclass[aspectratio=169, t, envcountsect, smaller, notheorems]{beamer}
\usetheme[sectionpage=simple, titleformat=allsmallcaps]{metropolis}
\usepackage[ngerman]{babel}
\usepackage{amssymb, mathtools, cancel, pgfplots, nicefrac, stmaryrd}
\usepackage[ruled, linesnumbered, vlined, ngerman, algosection]{algorithm2e}
\usepackage{array, csquotes, hyperref, tasks, multicol, wrapfig}

\title[Algorithmen und Datenstrukturen]{Algorithmen und Datenstrukturen}
\subtitle{Sommersemester 2019}
\author{David Luis Wiegandt}
\institute{Wissensmanagement in der Bioinformatik\\Institut für Informatik\\Humboldt-Universität zu Berlin}
\date{}
\titlegraphic{\includegraphics[width=70pt, height=70pt]{images/husiegel_bw.eps}}

\MakeOuterQuote{"}

\newcolumntype{L}[1]{>{\raggedleft\arraybackslash}p{#1}}

\usetikzlibrary{automata, graphs, positioning, arrows, patterns, fillbetween}

\tikzset{onslide/.code args={<#1>#2}{%
  \only<#1>{\pgfkeysalso{#2}}
}}

\setsansfont[Path = fonts/, UprightFont = *-Regular, BoldFont = *-Medium, ItalicFont = *-Italic]{FiraSans}

\resetcounteronoverlays{algocf}
\resetcounteronoverlays{theorem}

\settasks{label-format=\bfseries, counter-format=(tsk[a]), label-offset=.6em}

\uselanguage{german}
\languagepath{german}
\deftranslation[to=german]{Solution}{L\"osung}
\deftranslation[to=german]{Example}{Beispiel}
\deftranslation[to=german]{Task}{Aufgabe}
\deftranslation[to=german]{Remark}{Bemerkung}
\deftranslation[to=german]{Corollary}{Korollar}

\definecolor{hublue}{RGB}{0, 55, 108}
\definecolor{lighthublue}{RGB}{9, 80, 159}
\definecolor{tealblue}{RGB}{56, 134, 151}
\definecolor{madderlake}{RGB}{204, 41, 54}
\definecolor{ashgrey}{RGB}{174, 195, 176}
\definecolor{mandarin}{RGB}{239, 123, 69}
\definecolor{amaranth}{RGB}{160, 26, 125}

\newcommand<>{\emphlightblue}[1]{\textit#2{\color#2{lighthublue}#1}}
\newcommand<>{\emphteal}[1]{\textit#2{\color#2{teal}#1}}
\newcommand<>{\emphred}[1]{\textit#2{\color#2{madderlake}#1}}
\newcommand<>{\emphorange}[1]{\textit#2{\color#2{mandarin}#1}}
\newcommand<>{\emphpurple}[1]{\textit#2{\color#2{amaranth}#1}}

\newcommand<>{\lightblue}[1]{{\color#2{lighthublue}#1}}
\newcommand<>{\teal}[1]{{\color#2{teal}#1}}
\newcommand<>{\red}[1]{{\color#2{madderlake}#1}}
\newcommand<>{\orange}[1]{{\color#2{mandarin}#1}}
\newcommand<>{\purple}[1]{{\color#2{amaranth}#1}}

\makeatletter
\patchcmd{\beamer@sectionintoc}{\vfill}{\vskip\itemsep}{}{}

\setbeamertemplate{frame footer}{\insertshorttitle}
\setbeamertemplate{itemize item}{\fontsize{3}{4}\raise2.5pt\hbox{\donotcoloroutermaths$\blacksquare$}}
\setbeamertemplate{itemize subitem}{\fontsize{2.5}{3.5}\raise2.5pt\hbox{\donotcoloroutermaths$\blacksquare$}}
\setbeamertemplate{itemize/enumerate body begin}{\normalsize}
\setbeamertemplate{itemize/enumerate subbody begin}{\small}
\setbeamertemplate{section in toc}[sections numbered]
\setbeamertemplate{subsection in toc}[subsections numbered]
\setbeamertemplate{theorems}[numbered]

\setbeamercolor{alerted text}{fg=lighthublue}
\setbeamercolor{block title alerted}{fg=madderlake}
\setbeamercolor{footline}{fg=hublue}
\setbeamercolor{frametitle}{fg=hublue, bg=black!2}
\setbeamercolor{normal text}{fg=hublue}
\setbeamercolor{progress bar}{fg=hublue}
\setbeamercolor{title page header}{fg=white, bg=hublue}

\setbeamerfont{alerted text}{shape=\itshape}
\setbeamerfont{author}{series=\bfseries, size=\small}
\setbeamerfont{footline}{shape=\itshape}
\setbeamerfont{page number in head/foot}{size=\scriptsize}
\setbeamerfont{section title}{size=\LARGE}
\setbeamerfont{title}{size=\LARGE}

\setbeamertemplate{title page}{
\vspace*{-1pt}
\begin{beamercolorbox}[wd=\paperwidth,ht=0.618\paperheight]{title page header}
    \vfill
    \par
    \raggedright%
    \leftskip28pt%
    {\usebeamerfont{title} tutorium \\ \vspace{5pt} \inserttitle}
    \par
    \vspace*{3pt}
    {\usebeamerfont{subtitle}\insertsubtitle}
    \vspace*{-5pt}
    \vfill
    \tikz[remember picture, overlay] \node[anchor=center, draw=none] at (0.8\paperwidth,0) {\inserttitlegraphic};
\end{beamercolorbox}
\vspace*{2pt}
{\usebeamerfont{author}\beamer@shortauthor}

\medskip
\fontsize{10}{8}\selectfont
\par

{\usebeamerfont{institute}\insertinstitute}
\vfill
}
\makeatother

% default styles for plots
\pgfplotsset{width=\textwidth, compat=1.9}
\pgfplotscreateplotcyclelist{hubcols}{
    {tealblue},
    {madderlake},
    {ashgrey}
}
\pgfplotsset{
    tikzDefaults/.style={
        cycle list name=hubcols,
        font=\small,
        line width = 1pt,
        legend style={
            at={(0.5,-0.2)},
            anchor=north, 
            column sep=7pt,
            draw=none,
            fill=black!2
        },
        legend columns=2,
        axis x line=bottom,
        axis y line=middle,
        tick style={draw=hublue, semithick},
        tick align=center
        every axis/.append style={font=\tiny},
    },
    samples=100
}

\tikzset{
    node distance=2cm,
    every state/.style={
        thick,
        fill=lighthublue!7
    },
    every node/.style={
        draw,
        thick,
        circle
    },
    initial text={},
    double distance=2pt,
    every edge/.style={
        draw,
        ->,
        >=stealth,
        auto,
        thick,
        text=hublue,
        every node/.style={circle, fill=none, midway} %anchor=center
    }
}

% remove rules from algos while retaining the ruled layout
\setlength{\algoheightrule}{0pt}
\setlength{\algotitleheightrule}{0pt}

\let\leq=\leqslant
\let\geq=\geqslant

% general math
\newcommand{\nat}{\mathbb{N}}
\newcommand{\natpos}{\mathbb{N}_{\geq 1}}
\newcommand{\real}{\mathbb{R}}
\newcommand{\realpos}{\mathbb{R}_{\geq 0}}
\newcommand{\powerset}{\mathcal{P}}
\newcommand{\bigO}{\mathcal{O}}
\DeclareMathOperator{\cupdot}{\mathbin{\mathaccent\cdot\cup}}
\DeclareMathOperator{\subsetfin}{\mathbin{\subset_{\mathsf{e}}}}
\newcommand{\defiff}{\mathrel{\vcentcolon\Leftrightarrow}}
\DeclarePairedDelimiter\norm{\lVert}{\rVert}

% query functions
\DeclareMathOperator{\adom}{adom}
\DeclareMathOperator{\vars}{var}
\DeclareMathOperator{\cons}{adom}
\DeclareMathOperator{\free}{frei}
\newcommand{\qfn}[1]{[\![#1]\!]}

% algorithms
\SetKwInOut{Input}{input}
\SetKwInOut{Output}{output}
\ResetInOut{output}
\SetKw{True}{true}
\SetKw{False}{false}

% theorem environments
\theoremstyle{definition}
\newtheorem{theorem}[algocf]{\translate{Theorem}}
\newtheorem{lemma}[algocf]{\translate{Lemma}}
\newtheorem{definition}[algocf]{\translate{Definition}}
\newtheorem{example}[algocf]{\translate{Example}}
\newtheorem{task}[algocf]{\translate{Task}}
\newtheorem*{solution}{\translate{Solution}}
\newtheorem*{remark}{Bemerkung}

% enumerate environments
\newenvironment{subtasks}{
    \setbeamertemplate{enumerate item}{(\alph{enumi})}
    \setbeamerfont{enumerate item}{series=\bfseries}
    \begin{enumerate}
}{
    \end{enumerate}
}

\begin{document}
    \maketitle
    
    \begin{frame}{Organisatorisches}
    \begin{block}{Zeit \& Ort}
        \begin{itemize}
            \item Mittwoch von 15 bis 17 Uhr in RUD 25, 4.112
            \item Teilnahme ist \alert{freiwillig} und \alert{flexibel}
        \end{itemize}
    \end{block}
    
    \begin{block}{Inhalte}
        \begin{itemize}
            \item Fokus ist die \alert{Wiederholung} von Inhalten aus Vorlesung und Übung
            \item Dar\"uber hinaus ist auch die \alert{Vertiefung} von Inhalten aus der Vorlesung m\"oglich
        \end{itemize}
    \end{block}
    
    \begin{block}{Bei Fragen}
        \begin{itemize}
            \item Im Tutorium oder per Mail an \alert{\href{mailto:wigandtd@hu-berlin.de}{wigandtd@hu-berlin.de}}
        \end{itemize}
    \end{block}
    \end{frame}
    
    \begin{frame}{Behandelte Themen}
    \tableofcontents
    \end{frame}
    
    \section{Laufzeitanalyse}

\begin{frame}{Hintergrund}
\begin{block}{Knappheit von Ressourcen}
    \begin{itemize}
        \item Speicherplatz (Register, RAM, HDD/SSD, Tape, ...)
        \item \alert{Rechenleistung} (Taktfrequenz)
    \end{itemize}
\end{block}
\begin{block}{Eine Problemstellung -- viele Lösungsansätze}
    \begin{itemize}
        \item Welchen Ansatz sollen wir nutzen?
        \item Ziel: Klassifiziere Lösungsansätze nach ihrer \alert{Eignung}
        \begin{itemize}
            \item Eignung $\approx$ Ressourcenbedarf
        \end{itemize}
    \end{itemize}
\end{block}
\begin{block}{Berechnungsmodell}
    \begin{itemize}
        \item Wir analysieren \alert{abstrakte Algorithmen} -- folglich st\"utzen wir unsere Analysen auf \alert{abstrakte Maschinen}
    \end{itemize}
\end{block}
\end{frame}

\begin{frame}<1-2>[label=first_example]{Ein einfaches Beispiel}
\begin{example}
Gegeben sei eine Liste natürlicher Zahlen $a_1, \dots, a_n \in \nat$.

Maximiere $x = a_i \cdot a_j$ mit $1 \leq i < j \leq n$.

\begin{quote}
    In Worten: Finde das \alert{gr\"o{\ss}te} Produkt, das sich aus zwei \alert{verschiedenen} Elementen der Liste bilden lässt.
\end{quote}

\begin{block}<2->{M\"ogliche Ans\"atze}
    \begin{enumerate}
        \item Betrachte alle möglichen Paare und berechne deren Produkt
        \pause
        \item<4-> Suche die beiden gr\"o{\ss}ten Elemente $a_i, a_j$ und gib $a_i \cdot a_j$ aus
    \end{enumerate}
\end{block}
\end{example}
\end{frame}

\begin{frame}
\frametitle<1>{Algorithmus}
\frametitle<2>{Algorithmus \& Analyse}
\begin{columns}[T,onlytextwidth]
    \begin{column}{0.43\textwidth}
        \vspace{-5pt}
        \begin{algorithm}[H]
        	\caption{Maximales Produkt 1}
        	\label{alg:max_prod}
        	\DontPrintSemicolon
        	\Input{$a_1, \dots, a_n \in \nat$}
            \Output{$\max \{ a_i \cdot a_j : 1 \leq i < j \leq n \}$}
        	$x \gets 0$\;
        	\For{$i \gets 1$ \KwTo $n - 1$}{
        	    \For{$j \gets i + 1$ \KwTo $n$}{
        	        \If{$a_i \cdot a_j > x$}{
        	            $x \gets a_i \cdot a_j$\;
        	        }
                }
            }
            \Return{$x$}
        \end{algorithm}
    \end{column}
    \begin{column}<2>{0.51\textwidth}
        Die \"au{\ss}ere \textbf{for}-Schleife wird $(n-1)$-mal durchlaufen, die innere \textbf{for}-Schleife $n-1, n-2, \dots, 1$-mal
        \begin{itemize}
            \item Anwendung der \alert{gaußschen Summenformel} liefert $\frac{n \, (n-1)}{2}$
        \end{itemize}
        
        Der Aufwand des \textbf{if}-Statements ist von $n$ \alert{unabh\"angig}
        \begin{itemize}
            \item Multiplikationen, Vergleiche und Zuweisungen ben\"otigen lediglich \alert{konstant} viel Zeit
        \end{itemize}
        
        Insgesamt $\leq 4 \cdot \frac{n \, (n-1)}{2} + 2$ Instruktionen
    \end{column}
\end{columns}
\end{frame}

\againframe<3->{first_example}

\begin{frame}
\frametitle<1>{Algorithmus}
\frametitle<2>{Algorithmus \& Analyse}
\begin{columns}[T,onlytextwidth]
    \begin{column}{0.44\textwidth}
        \vspace{-5pt}
        \begin{algorithm}[H]
        	\caption{Maximales Produkt 2}
        	\label{alg:max_prod_ii}
        	\DontPrintSemicolon
        	\Input{$a_1, \dots, a_n \in \nat$}
            \Output{$\max \{ a_i \cdot a_j : 1 \leq i < j \leq n \}$}
        	$k \gets 0, \quad \ell \gets 0$\;
        	\For{$i \gets 1$ \KwTo $n$}{
        	    \If{$a_i > a_k$}{
        	        $\ell \gets k, \quad k \gets i$\;
    	        }
    	        \ElseIf{$a_i > a_{\ell}$}{
    	            $\ell \gets i$\;
    	        }
            }
            \Return{$a_k \cdot a_{\ell}$}
        \end{algorithm}
    \end{column}
    \begin{column}<2>{0.5\textwidth}
        Die \textbf{for}-Schleife wird genau $n$-mal durchlaufen
        
        \medskip
        
        Innerhalb der \textbf{if}-Statements f\"uhren wir lediglich grundlegende Operationen wie Vergleiche und Zuweisungen durch
        
        \medskip
        
        Insgesamt $\leq 3 \, n + 4$ Instruktionen
    \end{column}
\end{columns}
\end{frame}

\begin{frame}<1-2>[label=second_example]{Ein komplexeres Beispiel}
\begin{example}
Wir nennen eine Sequenz $a_1, \dots, a_n \in \nat$ \alert{partitionierbar}, falls  es $k, \ell \in \{1, \dots, n\}$ gibt, sodass $\sum_{i=k}^{\ell} a_i = \frac{1}{2} \cdot  \sum_{i=1}^{n} a_i$.

Entscheide f\"ur eine Sequenz $a_1, \dots, a_n$ mit $\{ a_1, \dots, a_n \} = \{1, \dots, n\}$, ob diese partitionierbar ist.

\begin{quote}
    In Worten: Entscheide f\"ur eine \alert{Permutation} der nat\"urlichen Zahlen von $1$ bis $n$, ob es eine \alert{zusammenh\"angende Subsequenz} gibt, deren Summe die \alert{H\"alfte} der Summe der Gesamtsequenz ist.
\end{quote}

\begin{block}<2->{M\"ogliche Ans\"atze}
    \begin{enumerate}
        \item Berechne die Summe aller möglichen Subsequenzen \emph{from scratch}
        \pause
        \item<4-> Nutze die Distributivit\"at der Addition, um redundante Berechnungen zu vermeiden
    \end{enumerate}
\end{block}
\end{example}
\end{frame}

\begin{frame}
\frametitle<1->{Algorithmus}
\frametitle<2->{Algorithmus \& Analyse}
\begin{columns}[T,onlytextwidth]
    \begin{column}{0.43\textwidth}
        \vspace{-5pt}
        \begin{algorithm}[H]
        	\caption{Subsequenz 1}
        	\label{alg:subseq_part_1}
        	\DontPrintSemicolon
        	\Input{$A = a_1, \dots, a_n \in \nat$}
            \Output{Ist $A$ partitionierbar?}
        	{\color<2>{madderlake}$s \gets \frac{n \, (n+1)}{4}$\;}
        	\For{$k \gets 1$ \KwTo $n$}{
        	    \For{$\ell \gets k$ \KwTo $n$}{
        	        {\color<3>{madderlake}$s' \gets 0$\;}
        	        \For{$i \gets k$ \KwTo $\ell$}{
        	            {\color<5>{madderlake}$s' \gets s' + a_i$\;}
                    }
                    \If{{\color<3>{madderlake}$s' = s$}}{
                        {\color<4>{madderlake}\Return{\True}}
                    }
                }
            }
            {\color<4>{madderlake}\Return{\False}}
        \end{algorithm}
    \end{column}
    \begin{column}<2->{0.51\textwidth}
        \only<2->{
            Die \emphred<2>{Initialisierung} ben\"otigt 4 Instruktionen
            \begin{itemize}
                \item Eine Addition, eine Multiplikation, eine Division und eine Zuweisung
            \end{itemize}
        }
        
        \only<3->{
            Die \emphred<3>{Zuweisung} in Zeile $4$ und der \emphred<3>{Vergleich} in Zeile $6$ werden h\"ochstens $\frac{n \, (n+1)}{2}$-mal ausgef\"uhrt
        }
        
        \medskip
        
        \only<4->{
            Einmalig wird ein Wert \emphred<4>{zurückgegeben}
        }
        
        \medskip
        
        \only<5->{
            Die \emphred<5>{innere \textbf{for}-Schleife} wird $\frac{n \, (n+1) \, (n+2)}{6}$-mal durchlaufen (Details: n\"achste Folie!)
            \begin{itemize}
                \item Entsprechend wird auch Zeile $6$ so oft ausgef\"uhrt
            \end{itemize}
        }
        
        \only<6->{
            Insgesamt $\leq \nicefrac{1}{3} \, n \, (n+1) \, (n+5) + 5$ Instruktionen
        }
    \end{column}
\end{columns}
\end{frame}

\begin{frame}{Details zur Laufzeitschranke}
    \begin{block}{Beobachtungen}
        \begin{itemize}
            \item<+-> Die innere \textbf{for}-Schleife wird stets $(\ell - k + 1)$-mal durchlaufen
            \begin{itemize}
                \item $\underbrace{\underbrace{1}_{\ell=1}, \underbrace{2}_{\ell=2}, \dots, \underbrace{n}_{\ell=n}}_{k=1}\text{-mal}, \underbrace{1, \dots, (n-1)}_{k=2}\text{-mal}, \dots, \underbrace{1, \dots, (n-n+1)}_{k=n}\text{-mal}$
            \end{itemize}
            \item<+-> Als Summe ergibt sich $$\sum_{i=1}^n \frac{i \, (i+1)}{2} = \frac{1}{2} \, \sum_{i=1}^n i \, (i+1) = \frac{1}{2} \, \left({\color<4>{madderlake} \sum_{i=1}^n i^2} + {\color<3->{mandarin} \sum_{i=1}^n i} \right)$$
            \item<+-> Wir wissen bereits, dass \emphorange{$\sum_{i=1}^n i = \frac{1}{2} \, n \, (n+1)$}
            \item<+-> Wir k\"onnen per vollst\"andiger Induktion zeigen, dass \emphred{$\sum_{i=1}^n i^2 = \frac{1}{6} \, n \, (n + 1) \, (2 \, n + 1)$}
        \end{itemize}
    \end{block}
\end{frame}

\begin{frame}{Details zur Laufzeitschranke (2)}
    \begin{lemma}\label{lem:nsquaressum}
    F\"ur alle $n \in \natpos$ gilt $\sum_{i=1}^n i^2 = \frac{1}{6} \, n \, (n + 1) \, (2 \, n + 1)$.
    \end{lemma}
    
    \begin{proof}
    Per Induktion.
    F\"ur $n=1$ gilt Lemma \ref{lem:nsquaressum} offensichtlich.
    F\"ur $n \leadsto n+1$ erhalten wir:
    \begin{align*}
        \sum_{i=1}^{n+1} i^2 = \sum_{i=1}^{n} i^2 + (n+1)^2 &\overset{!}{=} \frac{1}{6} \, n \, (n + 1) \, (2 \, n + 1) + (n+1)^2 \\
        &= \frac{1}{6} \, (n + 1) \, (n \, (2 \, n + 1) + 6 \, (n+1)) \\
        &= \frac{1}{6} \, (n + 1) \, (2 \, n^2 + 7 \, n + 6) \\
        &= \frac{1}{6} \, (n + 1) \, (n + 2) \, (2 \, n + 3)
    \end{align*}
    \end{proof}
\end{frame}

\begin{frame}{Details zur Laufzeitschranke (3)}
    Insgesamt erhalten wir f\"ur Algorithmus \ref{alg:subseq_part_1} unter Verwendung der \emphorange{gaußschen Summenformel} und \emphred{Lemma \ref{lem:nsquaressum}}:
    \begin{align*}
        \frac{1}{2} \, \left({\color{madderlake}\sum_{i=1}^n i^2} + {\color{mandarin}\sum_{i=1}^n i} \right) &= \frac{1}{2} \, \left({\color{madderlake}\frac{1}{6} \, n \, (n + 1) \, (2 \, n + 1)} + {\color{mandarin}\frac{1}{2} \, n \, (n+1)} \right) \\
        &= \frac{1}{2} \cdot \frac{1}{6} \, \left( n \, (n+1) \, (2\,n+1) + 3 \, n \, (n+1) \right) \\
        &= \frac{1}{2} \cdot \frac{1}{6} \, (2\,n+1+3) \, n \, (n+1) \\
        &= \frac{1}{2} \cdot \frac{1}{6} \cdot 2\, (n+2) \, n \, (n+1) \\
        &= \frac{1}{6} \, n \, (n+1) \, (n+2)
    \end{align*}
\end{frame}

\begin{frame}{Bemerkung zum Algorithmus}
    \begin{remark}
        In Algorithmus \ref{alg:subseq_part_1} summieren wir viele Subsequenzen mehrfach.
        
        Sei etwa $A = 4, 2, 3, 1$ mit $n=|A|=4$ die Eingabe.
        Zu Beginn ist $k=1$ und $\ell=1..4$, es werden also nacheinander die Summen $4$, $4+2$, $4+2+3$ und $4+2+3+1$ bestimmt.
        
        Tats\"achlich kann aber jede der Summen durch \alert{eine einzelne Addition} zum Ergebnis der vorherigen Summe bestimmt werden.
        
        So kann die Laufzeit von $\leq \nicefrac{1}{3} \, n \, (n+1) \, (n+5) + 5$ auf $\leq \nicefrac{1}{2} \, (3 \, n^2 + 5 \, n + 10)$ reduziert werden.
    \end{remark}
\end{frame}

\againframe<3->{second_example}

\begin{frame}{Algorithmus \& Analyse}
\begin{columns}[T,onlytextwidth]
\begin{column}{0.43\textwidth}
\vspace{-5pt}
\begin{algorithm}[H]
	\caption{Subsequenz 2}
	\label{alg:subseq_part_2}
	\DontPrintSemicolon
	\Input{$A = a_1, \dots, a_n \in \nat$}
    \Output{Ist $A$ partitionierbar?}
	$s \gets \frac{n \, (n+1)}{4}$\;
	\For{$k \gets 1$ \KwTo $n$}{
	    $s' \gets 0$\;
	    \For{$\ell \gets k$ \KwTo $n$}{
	        $s' \gets s' + a_i$\;
	        
            \If{$s' = s$}{
                \Return{\True}
            }
        }
    }
    \Return{\False}
\end{algorithm}
\end{column}
\begin{column}{0.51\textwidth}
    Die Initialisierung ben\"otigt 4 Instruktionen
    
    \medskip
    
    Der Rumpf der \"au{\ss}eren \textbf{for}-Schleife, und somit auch die Zuweisung in Zeile 3, wird h\"ochstens $n$-mal ausgef\"uhrt
    
    \medskip
    
    Der Rumpf der inneren \textbf{for}-Schleife, bestehend aus einer Addition, einer Zuweisung und einem Vergleich in den Zeilen 5 und 6, wird h\"ochstens $\frac{n \, (n+1)}{2}$-mal ausgef\"uhrt
    
    \medskip
    
    Einmalig wird ein Wert zur\"uckgegeben
    
    \medskip
    
    Insgesamt $\leq \nicefrac{1}{2} \, (3 \, n^2 + 5 \, n + 10)$ Instruktionen
\end{column}
\end{columns}
\end{frame}

\begin{frame}{Details zur Laufzeitschranke}
    Die angegebene Laufzeitschranke von $\nicefrac{1}{2} \, (3 \, n^2 + 5 \, n + 10)$ erhalten wir wie folgt:
    \begin{align*}
        4 + n + 3 \cdot \frac{n \, (n+1)}{2} + 1 &= \nicefrac{1}{2} \cdot 3 \, n \, (n+1) + n + 5 \\
        &= \nicefrac{1}{2} \cdot (3 \, n^2 + 3\, n) + n + 5 \\
        &= \nicefrac{1}{2} \cdot (3 \, n^2 + 3\, n) + \nicefrac{1}{2} \cdot (2 \, n + 10) \\
        &= \nicefrac{1}{2} \cdot (3 \, n^2 + 5\, n + 10)
    \end{align*}

    \begin{remark}
    Insbesondere haben wir damit gezeigt, dass \algref{alg:subseq_part_2} \alert{asymptotisch effizienter} ist als \algref{alg:subseq_part_1}, da die Laufzeit von \algref{alg:subseq_part_2} in $\bigO(n^2)$, die von  \algref{alg:subseq_part_1} hingegen in $\bigO(n^3)$ liegt (siehe Kapitel \nameref{sec:landau}).
    \end{remark}
\end{frame}

    \section{\uppercase{$\bigO$}-Notation}\label{sec:landau}

\begin{frame}{Definition}
\begin{columns}[T,onlytextwidth]
\begin{column}{0.47\textwidth}
    \begin{definition}[$\bigO$-Notation]
    Seien $\teal{f}, \red{g}: \nat \to \realpos$ Funktionen.
    Wir schreiben $\teal{f} \in \bigO(\red{g})$ genau dann, wenn es eine Konstante $\orange{c} \in \real_{> 0}$ gibt, sodass ab einem Index $\purple{n_0}$ f\"ur alle $n \geq \purple{n_0}$ gilt, dass $\teal{f}(n) \leq \orange{c} \cdot \red{g}(n)$.
    \end{definition}
    
    \medskip
    
    {\footnotesize Das hei{\ss}t, ab einem gewissen Punkt $\purple{n_0}$ gibt es keinen \alert{wesentlichen} Unterschied mehr zwischen den Funktionswerten von $\teal{f}$ und $\red{g}$.}
\end{column}
\begin{column}{0.47\textwidth}
    \begin{tikzpicture}[every node/.style={draw=none}]
        \begin{axis}[
            tikzDefaults, 
            ymax=280,
            xticklabels=\empty,
            yticklabels=\empty,
            legend columns=3,
            extra x ticks={6},
            extra x tick labels={$n_0$},
            extra x tick style={x tick label style={amaranth}, x tick style={amaranth}},
            set layers={background, main},
            clip mode=individual
        ]
            \addplot+[name path=f, domain=-4:12, draw=teal]{2*x^2 - 5*x + 30};
            \addlegendentry{$\teal{f}$}
            \addplot+[domain=-4:12, draw=madderlake]{x^2};
            \addlegendentry{$\red{g}$}
            \addplot+[draw=amaranth, dashed, line width=0.5pt, forget plot] coordinates {(6, 0) (6, 350)};
            \foreach \k in {1, ..., 8}{
                \addplot+[domain=-4:12, draw=mandarin!40, line width=0.5pt, forget plot, on layer=background]{1/2 * \k * x^2};
            }
            \addplot+[name path=gc, domain=-4:12, draw=mandarin, line width=0.6pt, on layer=background]{2 * x^2};
            \addlegendentry{$\orange{c} \cdot \red{g}$}
            %\addplot+[amaranth!17] fill between[of=f and gc, soft clip={domain=6:12}];
        \end{axis}
    \end{tikzpicture}
\end{column}
\end{columns}
\end{frame}

\begin{frame}{Beispiel}
\begin{example}
Sei \emphteal{$f(n) \coloneqq 3 \, n + 5$} und \emphred{$g(n) \coloneqq 2 \, n$}.

Um zu zeigen, dass $\teal{f} \in \bigO(\red{g})$ ist, w\"ahlen wir {$\orange{c=2}$} und erhalten
\begin{equation*}\label{complexity:ex:bigO1}
    \teal{\only<1>{3 \, n} \only<2->{\cancel{3 \, n}} + 5} \leq \orange{2} \cdot \red{2 \, n} = \only<1>{3 \, n} \only<2->{\cancel{3 \, n}} + n.
\end{equation*}
\only<2->{F\"ur $n \geq 5$ ist also $\teal{f(n)} \leq \orange{2} \cdot \red{g(n)}$ und somit $\purple{n_0 = 5}$.}
\end{example}

\begin{remark}<3->
In der Praxis schreibt man h\"aufig $f(n) \in \bigO(g(n))$ anstelle von $f \in \bigO(g)$.

Beispielsweise ist der Ausdruck $\teal{3 \, n + 5} \in \bigO(\red{2 \, n})$ \"aquivalent zu $$\teal{f} \in \bigO(\red{g}) \text{ f\"ur } \teal{f(n) \coloneqq 3 \, n + 5} \text{ und } \red{g(n) \coloneqq 2 \, n}.$$
\end{remark}
\end{frame}

\begin{frame}{Leichte Aufgaben}
\begin{task}<1->
Zeigen Sie, dass $6 \, n + 5 \in \bigO(n)$ ist.
\end{task}
\begin{solution}<2->
F\"ur $c = 7$ erhalten wir $6 \, n + 5 \leq 7 \, n = 6 \, n + n \iff n \geq 5$, also $n_0 = 5$.
\end{solution}

\begin{task}<3->
Zeigen Sie, dass gilt:
\begin{tasks}(3)
    \task $12 \, n - 4 \in \bigO(n)$
    \task $7 \, n + 1 \in \bigO(n^2)$
    \task $\log n^2 \in \bigO(\log n)$ 
\end{tasks}
\end{task}
\begin{solution}<4->
\begin{tasks}(3)
    \task Setze $c=12$ und $n_0=1$
    \task $7 \, n + 1 \leq n^2 \iff n \geq 8$, also $c=1$ und $n_0=8$
    \task Wegen $\log n^2 = 2 \, \log n$ setze $c=2$ und $n_0=1$
\end{tasks}
\end{solution}
\end{frame}

\begin{frame}{Wiederholung: Logarithmusgesetze}
\begin{theorem}
Seien $a, b \in \nat$ und sei $\log \coloneqq \log_x$ f\"ur ein beliebiges, aber festes $x \in \real_{>1}$ sowie $\ln \coloneqq \log_e$. Dann gilt:
\begin{align*}
    \log_b a &= \frac{\log_x a}{\log_x b}, \quad \text{ f\"ur alle $x \in \nat$}\tag{Basiswechsel} \\
    \log(a \cdot b) &= \log a + \log b\tag{Produktregel} \\
    \log(\nicefrac{a}{b}) &= \log a - \log b\tag{Quotientenregel} \\
    \log a^b &= b \, \log a\tag{Potenzregel}\\
    \derivative{n} \ln n &= \nicefrac{1}{n}\tag{Ableitungsregel}
\end{align*}
Falls nicht anders angegeben, nehmen wir zumeist $\log = \log_2$ an.
\end{theorem}
\end{frame}

\begin{frame}{Alternative Definition}
\begin{definition}<+->[$\bigO$-Notation]
Seien $\teal{f}, \red{g}: \nat \to \realpos$ Funktionen.
Wir schreiben $\teal{f} \in \bigO(\red{g})$ genau dann\footnote[frame]{Falls der Grenzwert \emph{nicht} existiert, nutzt man stattdessen den \alert{Limes Superior}: $\limsup_{n \to \infty} \frac{\teal{f(n)}}{\red{g(n)}} < \infty$.}, wenn $\red{g(n)} > 0$ f\"ur alle $n \in \nat$ und $\lim_{n \to \infty} \frac{\teal{f(n)}}{\red{g(n)}} < \infty$.
\end{definition}

\begin{example}<+->
F\"ur $f(n) \coloneqq 3 \, n + 5$ und $g(n) \coloneqq 2 \, n$ erhalten wir $$\lim_{n \to \infty} \frac{f(n)}{g(n)} = \lim_{n \to \infty} \frac{3 \, n + 5}{2 \, n} = \lim_{n \to \infty} \frac{3}{2} + \lim_{n \to \infty} \frac{5}{2 \, n} = \frac{3}{2} + 0 = \frac{3}{2} < \infty.$$
Folglich ist $3 \, n + 5 \in \bigO(2 \, n)$ --- was wir bereits aus Beispiel \ref{complexity:ex:bigO1} wissen.
\end{example}
\end{frame}

\begin{frame}<1>[label=hopital]{Der Satz von L'H\^opital}
\begin{theorem}[Satz von L'H\^opital]
F\"ur Funktionen $f, g : \nat \to \realpos$ mit $\lim_{n \to \infty} f(n) = \lim_{n \to \infty} g(n) \in \{ 0, \infty \}$ gilt: 
\begin{equation*}
    \lim_{n \to \infty} \frac{f(n)}{g(n)} = \lim_{n \to \infty} \frac{f'(n)}{g'(n)}.\only<3->{\tag{\red{*}}}
\end{equation*}
\end{theorem}

\pause

\begin{example}<3->
Sei $f(n) \coloneqq \ln n$ und $g(n) \coloneqq \sqrt{n}$.
Offensichtlich ist $\lim_{n \to \infty} \ln n = \lim_{n \to \infty} \sqrt{n} = \infty$.

Wir k\"onnen also den Satz von L'H\^opital anwenden und erhalten:
$$\lim_{n \to \infty} \frac{\ln n}{\sqrt{n}} \overset{(\red{\star})}{=} \lim_{n \to \infty} \frac{\frac{1}{n}}{\frac{1}{2 \, \sqrt{n}}} = \lim_{n \to \infty} \frac{2 \, \sqrt{n}}{n} = \lim_{n \to \infty} \frac{2}{\sqrt{n}} = 0 < \infty$$
Folglich ist $\ln n \in \bigO(\sqrt{n})$.
\end{example}
\end{frame}

\begin{frame}{Wiederholung: Ableitungsregeln I}
\begin{theorem}\label{complexity:thm:diff_rules}
Seien $f, g: \nat \to \realpos$ Funktionen und sei $(f \cdot g)(n) \coloneqq f(n) \cdot g(n)$.
Dann gilt:
\begin{align*}
    (f \cdot g)' &= f' \cdot g + f \cdot g' \tag{Produktregel} \\
    \left(\frac{f}{g}\right)' &= \frac{f' \cdot g - f \cdot g'}{g^2} \tag{Quotientenregel} \\
    (f \circ g)' &= (f' \circ g) \cdot g' \tag{Kettenregel}
\end{align*}
\end{theorem}

\begin{remark}
Die Quotientenregel folgt direkt aus der Produktregel und der Kettenregel: $$\left(\frac{f}{g}\right)' = (f \cdot g^{-1})' = f' \cdot g^{-1} + f \cdot (-1) \cdot g^{-2} \cdot g' = f' \cdot g \cdot g^{-2} - f \cdot g' \cdot g^{-2} = \frac{f' \cdot g - f \cdot g'}{g^2}$$
\end{remark}
\end{frame}

\begin{frame}{Wiederholung: Ableitungsregeln II}
\begin{task}
Sei $f_b(n) \coloneqq b^n$.
Zeigen Sie, dass f\"ur alle $b \in \natpos$ gilt: $$f_b'(n) = b^n \cdot \ln b$$
\end{task}

\begin{solution}<2->
Wir formen zun\"achst $f_b$ um und erhalten so: $$f_b(n) = \red{b^n = e^{\ln b^n}} = e^{n \, \ln b}.$$
Nun k\"onnen wir unter Verwendung der Kettenregel leicht $f_b'$ bestimmen:
$$f_b'(n) = \left(e^{n \, \ln b}\right)' = e^{n \, \ln b} \cdot \ln b = b^n \cdot \ln b$$
\end{solution}
\end{frame}

\begin{frame}{Wiederholung: Ableitungsregeln III}
\begin{task}
Leiten Sie folgende Funktionen ab, wobei $c \in \realpos$ konstant sei:
\begin{tasks}(3)
    \task $f(n) \coloneqq n^c \cdot \ln n$
    \task $f(n) \coloneqq \ln \ln n$
    \task\label{complexity:task:diff_nn} $f(n) \coloneqq n^n$
\end{tasks}
\end{task}

\begin{solution}<2->
\begin{tasks}
    \task $f'(n) = c \, n^{c-1} \cdot \ln n + n^c \cdot \nicefrac{1}{n} = n^{c-1} \, (c \cdot \ln n + 1)$
    \task $f'(n) = \frac{1}{\ln n} \cdot \nicefrac{1}{n} = \frac{1}{n \, \ln n}$
    \task Wir formen zun\"achst $f$ um und erhalten $f(n) = n^n = e^{\ln n^n} = e^{n \, \ln n}$.
    
    Nun bestimmen wir die Ableitung von $f$: 
    $$f'(n) = \left(e^{n \, \ln n}\right)' = \left(e^{n \, \ln n}\right) \cdot \left(1 \cdot \ln n + n \cdot \nicefrac{1}{n}\right) = n^n \cdot (\ln n + 1)$$
\end{tasks}
\end{solution}
\end{frame}

\againframe<2->{hopital}

\begin{frame}{Vermischte Aufgaben zur \MakeUppercase{$\bigO$}-Notation}
\begin{task}\label{complexity:task:various}
Ordnen Sie folgende Funktionen aufsteigend nach ihrem asymptotischen Wachstum und zeigen Sie die Korrektheit ihrer Anordnung: $$n!, \; 2^{\log n^3}, \; 2^n, \; \log(n^2 + 2^n), \; 2^{3 \, \log \frac{n}{2}}, \; n^n, \; \sqrt{n}$$
\end{task}

\begin{solution}<2->
Beobachte zun\"achst, dass $2^{\log n^3} = n^3$ und $2^{3 \, \log \frac{n}{2}} = 2^{\log \left( \left(\frac{n}{2}\right)^3\right)} = \left(\frac{n}{2}\right)^3 = \nicefrac{1}{8} \cdot n^3$.

Nun ist es leicht zu sehen, dass folgende Anordnung der Aufgabenstellung gen\"ugt: $$\sqrt{n}, \; \log(n^2 + 2^n), \; \underbrace{2^{3 \, \log \frac{n}{2}}}_{=\nicefrac{1}{8} \cdot n^3}, \; \underbrace{2^{\log n^3}}_{=n^3}, \; 2^n, \; n!, \; n^n$$
\end{solution}
\end{frame}

\begin{frame}{Details zur L\"osung}
\begin{itemize}
    \item Um $\sqrt{n} \in \bigO(\log(n^2 + 2^n))$ zu zeigen, sch\"atzen wir wie folgt ab: $$\sqrt{n} \leq n = \log 2^n \leq \log(n^2 + 2^n)$$
    Wir setzen also $c=n_0=1$ und sind fertig.
    \item Bzgl. $\log(n^2 + 2^n) \in \bigO(2^{3 \, \log \frac{n}{2}})$ setzen wir $c=8$ und $n_0=4$ und erhalten: $$\log(n^2 + 2^n) \leq \log(2^n + 2^n) = \log 2^{n+1} = n+1 \leq 8 \cdot \nicefrac{1}{8} \cdot n^3 = c \cdot 2^{3 \, \log \frac{n}{2}}$$
    \item $\nicefrac{1}{8} \cdot n^3 \in \bigO(n^3)$ gilt offensichtlich, da $\nicefrac{1}{8} \cdot n^3 \leq c \cdot n^3$ f\"ur alle $c \geq \nicefrac{1}{8}$ und $n \geq 1$.
    \item Wegen $\lim_{n \to \infty} \frac{n^3}{2^n} = \lim_{n \to \infty} \frac{3 \, n^2}{2^n \cdot \ln 2} = \lim_{n \to \infty} \frac{6 \, n}{2^n \cdot \ln^2 2} = \lim_{n \to \infty} \frac{6}{2^n \cdot \ln^3 2} = 0$ ist $n^3 \in \bigO(2^n)$
    \item F\"ur $2^n \in \bigO(n!)$ w\"ahlen wir $c=2$ und erhalten f\"ur alle $n \geq n_0 = 1$: $$2^n = \underbrace{2 \cdot 2 \cdot \cdots \cdot 2}_{n\text{ Terme}} \leq \underbrace{2 \cdot 2 \cdot 3 \cdot 4 \cdot \cdots \cdot n}_{n \text{ Terme}} = 2 \, n! = c \cdot n!$$
    \item Per Induktion nach $n$ folgt leicht, dass $n! \leq c \cdot n^n$ f\"ur beliebige $c \geq 1$ und alle $n \geq 1$
\end{itemize}
\end{frame}

    \section{Stacks \& Queues}

\begin{frame}{Abstrakte Datentypen}
\begin{block}{Grundlagen}
\begin{itemize}
    \item Ein \alert{abstrakter Datentyp (ADT)} beschreibt eine Menge von Operationen auf einer Menge von Werten
    \item Wir interessieren uns nur f\"ur die \alert{Signaturen} dieser Operationen
\end{itemize}
\end{block}

\begin{block}{Stacks \& Queues}
\begin{itemize}
    \item Eine \alert{Queue} (Warteschlange) ist ein ADT nach dem \alert{FIFO} (first-in, first-out) Prinzip
    \begin{itemize}
        \item Typische Operationen sind \texttt{enqueue} und \texttt{dequeue}
    \end{itemize}
    \item Ein \alert{Stack} (Stapel) ist ein ADT nach dem \alert{LIFO} (last-in, first-out) Prinzip
    \begin{itemize}
        \item Typische Operationen sind \texttt{push} und \texttt{pop}
    \end{itemize}
    \item Oft steht eine weitere Operation \texttt{head} zur Verf\"ugung, um das "n\"achste" Element zu betrachten, ohne den \alert{Zustand} der Datenstruktur zu modifizieren
\end{itemize}
\end{block}
\end{frame}

\begin{frame}[label=logexpr_intro]{Logische Ausdr\"ucke}
Im Folgenden f\"uhren wir \alert{Syntax} und \alert{Semantik} logischer Ausdr\"ucke ein
\begin{itemize}
    \item Die \alert{Syntax} beschreibt, \alert{wie} ein logischer Ausdruck aufgebaut sein muss
    \item Die \alert{Semantik} beschreibt, \alert{was} ein logischer Ausdruck aussagt
\end{itemize}

\begin{example}<2->
Wie wir gleich sehen werden, ist $\logic{(1 \land (1 \lor 0))}$ ein logischer Ausdruck.
\begin{itemize}
    \item Hierbei nehmen $\logic{1}$ und $\logic{0}$ die Rollen von "\emphteal{wahr}" und "\emphteal{falsch}" ein
    \item Um ihn "auszurechnen", m\"ussen wir u.a. festlegen, was das Ergebnis der Operatoren $\logic{\land}$ und $\logic{\lor}$ sein soll
    \item Intuitiv ist $\logic{(1 \land (1 \lor 0))} = \logic{(1 \land 1)} = \logic{1}$
\end{itemize}
\end{example}
\end{frame}

\begin{frame}{Syntax logischer Ausdr\"ucke}
\begin{definition}\label{sq:def:logexp_syn}
Wir definieren \emphteal{logische Ausdr\"ucke} als Worte \"uber dem Alphabet $$\Sigma = \{ \overbrace{\logic{0}, \, \logic{1}}^{\text{Ziffern}} \} \cup \{ \overbrace{\logic{\land}, \, \logic{\lor}, \, \logic{\oplus}}^{\text{Operatoren}}\} \cup \{ \overbrace{\logic{(}, \, \logic{)}}^{\text{Klammern}} \}.$$

Die Menge der logischen Ausdr\"ucke $\mathcal{L} \subseteq \Sigma^{\star}$ ist rekursiv wie folgt definiert:
\begin{itemize}
    \item Ziffern sind logische Ausdr\"ucke, d.h. $\{ \logic{0}, \logic{1} \} \subseteq \mathcal{L}$
    \item Falls $\teal{w_1}, \teal{w_2} \in \mathcal{L}$, so sind auch $\teal{(w_1 \land w_2)}, \teal{(w_1 \lor w_2)}, \teal{(w_1 \oplus w_2)} \in \mathcal{L}$
\end{itemize}
\end{definition}
\begin{example}<2->
Beispielsweise sind $\logic{1}$, $\logic{(0 \land 1)}$ und $\logic{(1 \land (1 \lor 0))}$ logische Ausdr\"ucke, w\"ahrend $\mathtt{\red{1 \land 1}}$ und $\mathtt{\red{(1 \land 1 \lor 0)}}$ \emph{keine} logischen Ausdr\"ucke sind.
\end{example}
\end{frame}

\begin{frame}{\"Uberpr\"ufung logischer Ausdr\"ucke}
\begin{task}
Sei $w \in \Sigma^{\star}$ mit $\Sigma = \{ \mathtt{0}, \mathtt{1} \} \cup \{ \mathtt{\land}, \mathtt{\lor}, \mathtt{\oplus} \} \cup \{ \mathtt{(}, \mathtt{)} \}$.
Entwerfen Sie einen Algorithmus, der bei Eingabe von $w$ in Zeit $\bigO(n)$ entscheidet, ob $w \in \mathcal{L}$ ist, wobei $n = |w|$ die L\"ange von $w$ in Zeichen sei.
\end{task}
\begin{solution}<2->
Wir arbeiten uns entsprechend der Klammerung von $w$ durch die Eingabe und nutzen den \alert{Call Stack}, um den Kontext vor Auswertung eines Teilausdrucks zu sichern
\begin{itemize}
    \item Dabei gehen wir gem\"a{\ss} der Syntax logischer Ausdr\"ucke vor, verlangen also, dass entweder $w \in \{ \mathtt{0}, \mathtt{1} \}$ oder $w = (w_1 \star w_2)$ mit $w_1, w_2 \in \mathcal{L}$ und $\star \in \{ \land, \lor, \oplus \}$
\end{itemize}
Beachte, dass wir $w$ dabei \emph{nicht} auswerten.
\end{solution}
\end{frame}

\begin{frame}{Algorithmus}
\begin{algorithm}[H]
	\caption{\"Uberpr\"ufung logischer Ausdr\"ucke}
	\label{sq:alg:check}
	\DontPrintSemicolon
    \SetKwFunction{check}{isExpression}
    \SetKwFunction{fsm}{fsm}
    \SetKwProg{Procedure}{procedure}{}{}
    \begin{multicols}{2}
    \Input{$w = w_1w_2\dots w_n \in \Sigma^{\star}$}
    \Output{$w \in \mathcal{L}$?}
    \BlankLine
    \Procedure{\check{$k$}}{
        \lIf{$w_k \in \{ \logic{0}, \logic{1} \}$}{
            \Return{$k$}
        }
        \lElseIf{$w_k = \logic{(}$}{
            \Return{\fsm{$k$}}
        }
        \lElse{\Return{$\bot$}}
    }
    \BlankLine
    \Return{\check{$1$} $= n$}\;
    \Procedure{\fsm{$k$}}{
        $s \gets 0$\;
        \For{$i \gets k + 1$ \KwTo $n$}{
            $\ell \gets$ \check{$i$}\;
            \If{$s \in \{ 0, 2\}$ \And $\ell > 0$}{
                $s \gets s + 1, \quad i \gets \ell$\;
            }
            \ElseIf{$s = 1$ \And $w_i \in \{ \logic{\land}, \logic{\lor}, \logic{\oplus} \}$}{
                $s \gets s + 1$\;
            }
            \ElseIf{$s = 3$ \And $w_i = \logic{)}$}{
                \Return{$i$}\;
            }
            \lElse{\Return{$\bot$}}
        }
    }
    \end{multicols}
\end{algorithm}
\end{frame}

% \begin{frame}{Analyse}
% \end{frame}

\begin{frame}{Alternativer Ansatz \, I}
\begin{block}{Grundlagen}
Wir k\"onnen die Sprache aller logischen Ausdr\"ucke $\mathcal{L}$ leicht durch Angabe einer \alert{kontextfreien Grammatik} spezifizieren
\begin{itemize}
    \item Dementsprechend k\"onnen wir einen \alert{Kellerautomaten} (PDA) angeben, der f\"ur Worte $w \in \Sigma^{\star}$ entscheidet, ob $w \in \mathcal{L}$
    \begin{itemize}
        \item Zu Details siehe Vorlesung \emph{Einf\"uhrung in die Theoretische Informatik}
    \end{itemize}
\end{itemize}
\end{block}

\begin{block}{Idee}<2->
Wir simulieren den Lauf eines PDA algorithmisch
\begin{itemize}
    \item Um einen PDA zu simulieren, ben\"otigen wir einen \alert{Stack}
    \item Die \alert{\"Uberf\"uhrungsfunktion} implementieren wir durch eine \alert{Lookup-Tabelle}
\end{itemize}

Nun k\"onnen wir jedes Zeichen der Eingabe in Zeit $\bigO(1)$ verarbeiten --- insgesamt $\bigO(n)$
\end{block}
\end{frame}

\begin{frame}[fragile]{Alternativer Ansatz \, II}
\begin{example}
{
\newcommand{\transition}[3]{\logic{#1}\purple{#2}{\color{hublue},}\,\purple{#3}}
\newcommand{\hash}{\text{\#}}
Folgender PDA $M = (Z, \teal{\Sigma}, \purple{\Gamma}, \delta, q_0, \purple{\text{\#}})$ mit $Z=\{q_0, q_1, q_2, q_3, q_4\}$ und $\purple{\Gamma = \{ \text{\#}, L, R \}}$ erkennt genau die logischen Ausdr\"ucke gem\"a{\ss} Definition \ref{sq:def:logexp_syn}:

\begin{figure}[!h]
\centering
\begin{tikzpicture}[node distance=2.2cm, every loop/.style={looseness=5, out=115, in=65}]
    \node[state, initial] (q0) {$q_0$};
    \node[state, right of=q0] (q1) {$q_1$};
    \node[state, right of=q1] (q2) {$q_2$};
    \node[state, right of=q2] (q3) {$q_3$};
    \node[state, right of=q3] (q4) {$q_4$};
    \draw   (q0) edge[loop above] node[align=center]{$\transition{0}{\hash}{\epsilon}$ \\ $\transition{1}{\hash}{\epsilon}$} (q0)
            (q0) edge[above] node{$\transition{(}{\hash}{L}$} (q1)
            (q1) edge[loop above] node{$\transition{(}{L}{LL}$} (q1)
            (q1) edge[above] node[align=center]{$\logic{0}$ \\ $\logic{1}$} (q2)
            (q2) edge[above] node[align=center]{$\transition{\land}{L}{R}$ \\ $\transition{\lor}{L}{R}$ \\ $\transition{\oplus}{L}{R}$} (q3)
            (q3) edge[bend left, above] node[align=center]{$\logic{0}$ \\ $\logic{1}$} (q4)
            (q3) edge[bend left, below] node{$\transition{(}{R}{RL}$} (q1)
            (q4) edge[bend left, below] node[align=center]{$\transition{\land}{L}{R}$ \\ $\transition{\lor}{L}{R}$ \\ $\transition{\oplus}{L}{R}$} (q3)
            (q4) edge[loop above] node{$\transition{)}{R}{\epsilon}$} (q4);
\end{tikzpicture}
\caption{Ein PDA $M$ mit $L(M) = \mathcal{L}$}
\label{sq:img:pda}
\end{figure}
}
\end{example}
\end{frame}

\againframe<3->{logexpr_intro}

\begin{frame}{Semantik logischer Ausdr\"ucke}
\begin{definition}
Wir definieren den \emph{Wert} $\llbracket \teal{w} \rrbracket$ logischer Ausdr\"ucke $\teal{w} \in \mathcal{L}$ rekursiv wie folgt:
\begin{itemize}
    \item $\eval{\logic{0}} = \logic{0}$ und $\eval{\logic{1}} = \logic{1}$
    \item F\"ur logische Ausdr\"ucke $\teal{w_1}, \teal{w_2}$ ist
    \begin{itemize}
        \item $\eval{\teal{(w_1 \land w_2)}} = \eval{\teal{w_1}} \cdot \eval{\teal{w_2}}$
        \item $\eval{\teal{(w_1 \lor w_2)}} = \begin{cases}\logic{1}, & \text{falls } \eval{\teal{w_1}} + \eval{\teal{w_2}} \geq 1 \\ \logic{0}, & \text{sonst.}\end{cases}$
        \item $\eval{\teal{(w_1 \oplus w_2)}} = \begin{cases}\logic{1}, & \text{falls } \eval{\teal{w_1}} \neq \eval{\teal{w_2}} \\ \logic{0}, & \text{sonst.}\end{cases}$
    \end{itemize}
\end{itemize}
\end{definition}
\begin{example}<2->\label{sq:ex:sem}\vspace*{-6pt}
Sei $\logic{((0 \lor 1) \land (1 \land 0))}$ gegeben.
Wir erhalten $\eval{
\logic{(}
\underbrace{\logic{(0 \lor 1)}}_{\eval{\logic{(0 \lor 1)}} = \logic{1}}
\logic{\land}
\underbrace{\logic{(1 \land 0)}}_{\eval{\logic{(1 \land 0)}} = \logic{0}}
\logic{)}
} = \logic{0}$.
\end{example}
\end{frame}

\begin{frame}{Auswertung logischer Ausdr\"ucke}
\begin{task}
Entwerfen Sie einen Algorithmus, der einen logischen Ausdruck $\teal{w} \in \mathcal{L}$ auswertet, also $\eval{\teal{w}}$ bestimmt.
Die Laufzeit soll in $\bigO(n)$ liegen, wobei $n = |\teal{w}|$ die L\"ange des logischen Ausdrucks in Zeichen sei.

Nebst einer konstanten Anzahl zus\"atzlicher Variablen primitiven Datentyps d\"urfen Sie eine Queue oder einen Stack benutzen.
\end{task}
\begin{solution}<2->
\begin{quote}
    Idee: Wir werten \teal{w} wie in Beispiel \ref{sq:ex:sem} ``von innen nach au{\ss}en'' aus und merken uns Zwischenergebnisse in einem \alert{Stack}.
\end{quote}

Details folgen!
\end{solution}
\end{frame}
    \section{Sortierverfahren}

\begin{frame}{Grundlagen}
\begin{block}{Listen}
\begin{itemize}
    \item Abstrakter Datentyp zur Repr\"asentation geordneter Sequenzen von Datenwerten
    \item Verschiedene Implementierungen m\"oglich, bspw. als \alert{Array} oder durch \alert{Verkettung}
\end{itemize}
\end{block}

\begin{block}{Vereinbarungen}
Unabh\"angig von der konkreten Implementierung nehmen wir stets an, dass
\begin{itemize}
    \item Listen eine endliche L\"ange $n \in \nat$ haben
    \item Elemente einer Liste durch die Positionen $1, \dots, n$ adressiert werden
    %\item Operationen \texttt{insert}, \texttt{delete} und \texttt{elementAt} zur Verf\"ugung stehen
\end{itemize}
Den Zugriff auf das $i$-te Element einer Liste $L$ notieren wir mit $L_i$.
Die Teilliste vom $i$-ten bis zum $j$-ten Element (mit $i \leq j$) bezeichnen wir mit $L[i, j]$.
\end{block}
\end{frame}

\begin{frame}<1-3>[label=general_sort]
\frametitle<1-3>{Allgemeine Sortierverfahren}
\frametitle<4>{Zur Erinnerung: Allgemeine Sortierverfahren}
\begin{definition}[Allgemeines Sortierverfahren]
Sei $L = \alglist{a_1, \dots, a_n}$ eine Liste \"uber dem durch $\leq_D$ linear geordneten Definitionsbereich $D$.

Ein \alert{\red<2>{allgemeines} Sortierverfahren} berechnet eine Liste $L' = \alglist{a_{\pi(1)}, \dots, a_{\pi(n)}}$ f\"ur eine Permutation $\pi$ von $\{ 1, \dots, n\}$, sodass
\begin{enumerate}
    \item $a_{\pi(i)} \leq_D a_{\pi(i+1)}$ f\"ur alle $1 \leq i < n$ und
    \item<2-> \emphred<2>{nebst $\leq_D$ keine weiteren evtl. auf $D$ definierten Operationen genutzt werden}.
\end{enumerate}
\end{definition}

\begin{remark}<3->
Obgleich wir \emph{oft} Listen \"uber den nat\"urlichen Zahlen $\nat$ betrachten, d\"urfen wir \emph{keinen} Gebrauch von Operationen wie etwa Multiplikation oder Test auf \alert{Kongruenz} machen.
\end{remark}
\pause
\end{frame}

\begin{frame}{\"Uberblick \"uber allgemeine Sortierverfahren}
In der Vorlesung wurden unter anderem folgende \emph{allgemeine} Sortierverfahren vorgestellt:
\begin{itemize}
    \item Bubble Sort
    \item Insertion Sort
    \item Merge Sort
    \item Quick Sort
\end{itemize}
Im Folgenden f\"uhren wir \emph{Schreibtischtests} o.g. Verfahren durch.
\end{frame}

\begin{frame}{Bubble Sort}
\begin{columns}[T,onlytextwidth]
\begin{column}{0.4\textwidth}
\begin{algorithm}[H]
	\caption{Bubble Sort}
	\label{sort:alg:bubblesort}
	\DontPrintSemicolon
	\Input{Liste $L$}
    \Output{Sortierte Liste $L'$}
	\For{$i \gets 1$ \KwTo $n$}{
	    \For{$j \gets 1$ \KwTo $n - i$}{
            \If{$L_j \geq L_{j+1}$}{
                $t \gets L_j$\;
                $L_j \gets L_{j+1}$\;
                $L_{j+1} \gets t$\;
            }
        }
    }
    \Return{$L$}\;
\end{algorithm}
\end{column}
\begin{column}{0.55\textwidth}
\begin{task}<2->
Ordnen Sie folgende Liste mittels des nebenstehenden Algorithmus: $$L = \alglist{17, 5, 4, 63, 46, 58, 4}$$
\end{task}
% \begin{solution}<3->

% \end{solution}
\end{column}
\end{columns}
\end{frame}

\begin{frame}{Insertion Sort}
\begin{columns}[T,onlytextwidth]
\begin{column}{0.4\textwidth}
\begin{algorithm}[H]
	\caption{Insertion Sort}
	\label{sort:alg:insertionsort}
	\DontPrintSemicolon
	\Input{Liste $L$}
    \Output{Sortierte Liste $L'$}
	\For{$i \gets 2$ \KwTo $n$}{
	    $key \gets L_i$\;
	    $j \gets i - 1$\;
	    \While{$L_j > key$ \And $j > 0$}{
            $L_{j+1} \gets L_j$\;
            $j \gets j - 1$\;
        }
    }
    \Return{$L$}\;
\end{algorithm}
\end{column}
\begin{column}{0.55\textwidth}
\begin{task}<2->
Ordnen Sie folgende Liste mittels des nebenstehenden Algorithmus: $$L = \alglist{17, 5, 4, 63, 46, 58, 4}$$
\end{task}
% \begin{solution}<3->

% \end{solution}
\end{column}
\end{columns}
\end{frame}

\begin{frame}{Mergen sortierter Arrays}
Seien $A = \alglist{a_1, \dots, a_m}$ und $B = \alglist{a_{m+1}, \dots, a_{m+n}}$ sortierte Listen der L\"ange $m$ bzw. $n$.
Wir wollen nun eine ebenfalls sortierte Liste $L = \alglist{ a_{\pi(1)}, \dots, a_{\pi(m + n)} }$ bestimmen.
\begin{itemize}
    \item Das geht nat\"urlich, indem wir die Listen $A$ und $B$ zu $L' = \alglist{a_1, \dots, a_m, a_{m+1}, \dots, a_{m+n}}$ konkatenieren und $L'$ dann in $\bigO((m + n) \cdot \log(m + n))$ sortieren
    \item Damit nutzen wir aber nicht, dass $A$ und $B$ bereits sortiert sind
\end{itemize}
Stattdessen durchlaufen wir beide Listen "abwechselnd" von vorn nach hinten und f\"ugen den jeweils kleineren Wert ans Ende eines neuen Arrays ein
\begin{itemize}
    \item Nachteil: Wir ben\"otigen $\bigO(m+n)$ zus\"atzlichen Speicherplatz
\end{itemize}
\end{frame}

\begin{frame}{Mergen sortierter Arrays --- Beispiel}
\begin{example}
Sei $A = \alglist{ 2, 3, 5, 7, 11, 12 }$ und $B = \alglist{ 4, 11, 16, 17 }$.

\medskip

\begin{tabular}[t]{lr}
$A$: &
\begin{tabular}[t]{r|L{0.4cm}|}
    \cline{2-2} 
    \only<1>{$\rightarrow$} & 2 \\ \cline{2-2}
    \only<2>{$\rightarrow$} & 3 \\ \cline{2-2}
    \only<3-4>{$\rightarrow$} & 5 \\ \cline{2-2}
    \only<5>{$\rightarrow$} & 7 \\ \cline{2-2}
    \only<6>{$\rightarrow$} & 11 \\ \cline{2-2}
    \only<7-11>{$\rightarrow$} & 12 \\ \cline{2-2}
\end{tabular}
\end{tabular}
\hfill
\begin{tabular}[t]{lr}
$B$: &
\begin{tabular}[t]{r|L{0.4cm}|}
    \cline{2-2} 
    \only<1-3>{$\rightarrow$} & 4 \\ \cline{2-2}
    \only<4-7>{$\rightarrow$} & 11 \\ \cline{2-2}
    \only<8-9>{$\rightarrow$} & 16 \\ \cline{2-2}
    \only<10-11>{$\rightarrow$} & 17 \\ \cline{2-2}
\end{tabular}
\end{tabular}
\hfill
\begin{tabular}[t]{lr}
$L$: &
\begin{tabular}[t]{r|L{0.4cm}|}
    \cline{2-2} 
    \only<1>{$\rightarrow$} & \only<2->{2} \\ \cline{2-2}
    \only<2>{$\rightarrow$} & \only<3->{3} \\ \cline{2-2}
    \only<3>{$\rightarrow$} & \only<4->{4} \\ \cline{2-2}
    \only<4>{$\rightarrow$} & \only<5->{5} \\ \cline{2-2}
    \only<5>{$\rightarrow$} & \only<6->{7} \\ \cline{2-2}
    \only<6>{$\rightarrow$} & \only<7->{11} \\ \cline{2-2}
    \only<7>{$\rightarrow$} & \only<8->{11} \\ \cline{2-2}
    \only<8>{$\rightarrow$} & \only<9->{12} \\ \cline{2-2}
    \only<9>{$\rightarrow$} & \only<10->{16} \\ \cline{2-2}
    \only<10-11>{$\rightarrow$} & \only<11->{17} \\ \cline{2-2}
\end{tabular}
\end{tabular}
\end{example}
\end{frame}

\begin{frame}{Mergen sortierter Arrays --- Algorithmus}
\begin{algorithm}[H]
	\caption{Merge}
	\label{sort:alg:merge}
	\DontPrintSemicolon
	\Input{Sortierte Listen $A = \alglist{a_1, \dots, a_m}$ und $B = \alglist{a_{m+1}, \dots, a_{m+n}}$}
    \Output{Sortierte Liste $L = \alglist{a_{\pi(1)}, \dots, a_{\pi(m+n)}}$}
    \BlankLine
    \begin{multicols}{2}
    $i \gets 1, \quad j \gets m + 1, \quad k \gets 1$\;
    \While{$i \leq m$ \And $j \leq m + n$}{
        \If{$A_i \leq B_j$}{
            $L_k \gets A_i, \quad i \gets i + 1$
        }
        \Else{
            $L_k \gets B_j, \quad j \gets j + 1$
        }
        $k \gets k + 1$\;
    }
    \If{$i > m$}{
        $L[k, m+n] \gets B[j, m+n]$
    }
    \Else{
        $L[k, m+n] \gets A[i, m]$
    }
    \Return{$L$}\;
    \end{multicols}
\end{algorithm}
\end{frame}

\begin{frame}{Merge Sort}
\begin{columns}[T,onlytextwidth]
\begin{column}{0.42\textwidth}
\begin{algorithm}[H]
	\caption{Merge Sort}
	\label{sort:alg:mergesort}
	\DontPrintSemicolon
	\Input{Liste $L$}
    \Output{Sortierte Liste $L'$}
    \SetKwFunction{MergeSort}{MergeSort}
    \SetKwFunction{Merge}{Merge}
    \SetKwProg{Procedure}{procedure}{}{}
    \Procedure{\MergeSort{$L$}}{
        $n \gets |L|, \quad m \gets \floor{\nicefrac{n}{2}}$\;
        \BlankLine
        \lIf{$n \leq 1$}{
            \Return{$L$}
        }
        \BlankLine
        $L_{\ell} \gets$ \MergeSort{$L[1, m]$}\;
        $L_r \gets$ \MergeSort{$L[m + 1, n]$}\;
        \BlankLine
        \Return{\Merge{$L_{\ell}$, $L_r$}}\;
    }
    \Return{\MergeSort{$L$}}\;
\end{algorithm}
\end{column}
\begin{column}{0.53\textwidth}
\begin{task}<2->
Ordnen Sie folgende Liste mittels des nebenstehenden Algorithmus: $$L = \alglist{17, 5, 4, 63, 46, 58, 4}$$
\end{task}
% \begin{solution}<3->

% \end{solution}
\end{column}
\end{columns}
\end{frame}

\begin{frame}{Partitionierung von Listen}
\begin{definition}[Partitionierung]
Wir nennen eine Liste $L = \alglist{a_1, \dots, a_n}$ \alert{partitioniert bzgl. $a_k$} f\"ur ein $k \in \{ 1, \dots, n \}$, falls $a_i \leq a_k$ f\"ur alle $i < k$ und $a_j \geq a_k$ f\"ur alle $j > k$:
$$L = \alglist{\underbrace{a_1, \dots, a_{k-1}}_{\leq a_k}, a_k, \underbrace{a_{k+1}, \dots, a_n}_{\geq a_k}}$$
\end{definition}

\begin{example}<2->
Die Liste $L = \alglist{4, 7, 3, \teal{8}, 12, 9, \red{8}, 11}$ ist partitioniert bzgl. $\teal{8}$, \emph{nicht} aber bzgl. $\red{8}$.
\end{example}
\end{frame}

\begin{frame}{Partitionierung von Listen --- Problemstellung}
\begin{block}{Problemstellung}
Gegeben eine Liste $L = \alglist{a_1, \dots, a_n}$ und ein $k \in \{ 1, \dots, n \}$.
Bestimme in Zeit $\bigO(n)$ eine Liste $L' = \alglist{a_{\pi(1)}, \dots, a_{\pi(n)}}$, die bzgl. $a_{\pi^{-1}(k)}$ partitioniert ist.
\end{block}

\begin{remark}<2->
Da wir sp\"ater auf den Teillisten $L'[1, \pi^{-1}(k)-1]$ und $L'[\pi^{-1}(k)+1, n]$ links und rechts von Position $\pi^{-1}(k)$ rekursiv absteigen werden, ist idealerweise $\pi^{-1}(k) \approx \nicefrac{n}{2}$.
\end{remark}

\begin{block}{Frage}<3->
Wie k\"onnen wir ein solches $k$ in Zeit h\"ochstens $\bigO(n)$ finden?
\end{block}

\begin{block}{Antwort}<4->
Das ist gleichbedeutend damit, den \alert{Median} von $L$ in Zeit $\bigO(n)$ zu bestimmen --- das ist m\"oglich, aber nicht trivial (siehe bspw. \href{https://www.cs.hhu.de/fileadmin/redaktion/Fakultaeten/Mathematisch-Naturwissenschaftliche_Fakultaet/Informatik/Computational_Cell_Biology/AlDat/Alg2016_03.pdf}{\alert{hier}})

Eine zuver\"assige \alert{Heuristik} ist es, einen ``zuf\"alligen'' Wert zu w\"ahlen --- bspw. immer $a_n$
\end{block}
\end{frame}

\begin{frame}{Partitionierung von Listen --- Algorithmus}
\begin{algorithm}[H]
	\caption{Partitionierung}
	\label{sort:alg:partition}
	\DontPrintSemicolon
	\SetKwFunction{Swap}{swap}
	\SetKwRepeat{Do}{do}{while}
	\Input{Liste $L = \alglist{a_1, \dots, a_n}$}
    \Output{Bzgl. $a_n$ partitionierte Liste $L' = \alglist{a_{\pi(1)}, \dots, a_{\pi(n)}}$}
    $i \gets 1, \quad j \gets n - 1$\;
    \Do{$i < j$}{
        \lWhile{$L_i \leq L_n$ \And $i < n$}{$i \gets i + 1$}
        \lWhile{$L_j \geq L_n$ \And $j > 1$}{$j \gets j - 1$}
        \lIf{$i < j$}{\Swap{$L_i$, $L_j$}}
    }
    \BlankLine
    \Swap{$L_i$, $L_n$}\;
    \BlankLine
    \Return{$L[1, i - 1], L[i + 1, n]$}\;
\end{algorithm}
\end{frame}

\begin{frame}{Quick Sort}
\begin{columns}[T,onlytextwidth]
\begin{column}{0.42\textwidth}
\begin{algorithm}[H]
	\caption{Quick Sort}
	\label{sort:alg:quicksort}
	\DontPrintSemicolon
	\Input{Liste $L$}
    \Output{Sortierte Liste $L'$}
    \SetKwFunction{QuickSort}{QuickSort}
    \SetKwFunction{Merge}{Merge}
    \SetKwProg{Procedure}{procedure}{}{}
    \Procedure{\QuickSort{$L$}}{
        \lIf{$|L| \leq 1$}{\Return{$L$}}
        $L^L, L^R \gets \text{partitioniere } L \text{ bzgl. } L_n$\;
        $L' \gets $ \QuickSort{$L^L$}\;
        $R' \gets $ \QuickSort{$L^R$}\;
    }
    \Return{$\alglist{L', L_n, R'}$}\;
\end{algorithm}
\end{column}
\begin{column}{0.53\textwidth}
\begin{task}<2->
Ordnen Sie folgende Liste mittels des nebenstehenden Algorithmus: $$L = \alglist{17, 5, 4, 63, 46, 58, 4}$$
\end{task}
% \begin{solution}<3->

% \end{solution}
\end{column}
\end{columns}
\end{frame}

% \begin{frame}{Problemreduktion}
% Problemreduktionen sind ein in der theoretischen Informatik h\"aufig genutztes Werkzeug, um per \alert{Widerspruch} zu zeigen, dass ein Problem \emph{nicht} in einer bestimmten Komplexit\"atsklasse liegt.

% \begin{block}{Grundidee}<2->
% Sei $a$ eine Eingabe f\"ur ein Problem $\mathcal{A}$ und analog $b$ eine Eingabe f\"ur ein Problem $\mathcal{B}$.

% Wenn wir 
% \begin{enumerate}
%     \item einen Algorithmus $\mathbb{B}$ haben, der eine Eingabe $b$ f\"ur $\mathcal{B}$ verarbeiten kann (kurz: $\mathbb{B}(b)$ \emph{und} 
%     \item einen Algorithmus $\mathbb{T}$, der eine Eingabe $a$ f\"ur $\mathcal{A}$ in eine Eingabe f\"ur $\mathcal{B}$ \alert{transformieren} kann,
% \end{enumerate}
% dann k\"onnen wir $b \coloneqq \mathbb{F}(a)$ als Eingabe f\"ur $\mathbb{B}$ nutzen und so 

% \end{block}

% \begin{example}<3->
% Siehe Tafel!
% \end{example}
% \end{frame}

\begin{frame}{Bemerkung zur Laufzeit}
In der Vorlesung wurde folgendes Theorem bewiesen:

\begin{theorem}\label{sort:thm:bound}
Sei $L = \alglist{a_1, \dots, a_n}$ eine Liste mit $a_i \neq a_j$ f\"ur $1 \leq i < j \leq n$.
Um $L$ zu sortieren, ben\"otigt jedes \emph{allgemeine Sortierverfahren} im \emph{Worst Case} $\Omega(n \, \log n)$ Vergleiche.
\end{theorem}

\begin{block}{Bedeutung}
Das bedeutet insbesondere, dass wir f\"ur jeden der vorgestellten Algorithmen Instanzen erzeugen k\"onnen, f\"ur deren Sortierung mindestens $n \, \log n$ Vergleiche ben\"otigt werden.

Das gilt aber nur, wenn wir von \emph{beliebigen} Listen als Eingabe ausgehen.
\begin{itemize}
    \item Andernfalls k\"onnen wir gelegentlich durch zus\"atzliche \emph{Annahmen \"uber die Beschaffenheit} der Liste die untere Schranke f\"ur \emph{beliebige} Listen unterbieten.
\end{itemize}
\end{block}
\end{frame}

\begin{frame}{Sortieren spezieller Arrays}
Gelegentlich besitzen wir zus\"atzliches Wissen \"uber die Beschaffenheit einer eingegebenen Liste $L$.
Beispiele daf\"ur sind etwa:
\begin{itemize}
    \item "$L$ ist so vorsortiert, dass alle Elemente der ersten Hälfte von $L$ kleiner sind als alle Elemente der zweiten Hälfte."
    \item "Die Länge von $L$ ist durch 10 teilbar, und jeweils 10 aufeinanderfolgende Elemente (vom Anfang beginnend) sind aufsteigend sortiert."
\end{itemize}
Hierbei stellen wir uns die Frage, inwiefern sich diese Eigenschaften ggf. zur Reduktion der Laufzeit f\"ur's Sortieren nutzen lassen --- etwa auf $\bigO(n)$.

\begin{alertblock}{Achtung}<2->
Die einzelnen Datenwerte d\"urfen wir aber nach wie vor \emphred{nur miteinander vergleichen}!
\end{alertblock}
\end{frame}

% \begin{frame}{Entscheidungsprobleme}
% \begin{definition}[Entscheidungsproblem]
% Sei $U$ eine Menge und sei $A \subseteq U$.
% Das Problem, für ein $x \in U$ zu \emph{entscheiden}, ob auch $x \in A$ ist, wird das \alert{Entscheidungsproblem} für $A$ genannt.

% Ein Algorithmus $\mathbb{A}: U \to \{ 0, 1 \}$ \alert{entscheidet} $A$, falls $\mathbb{A}(x) = 1 \iff x \in A$.

% Das Entscheidungsproblem f\"ur $A$ liegt in $\bigO(g)$ für eine Funktion $g$, falls ein Algorithmus existiert, der $A$ in Zeit $\bigO(g)$ entscheidet.
% \end{definition}

% \begin{example}<2->
% Sei $U = \nat$ die Menge \emph{aller} nat\"urlichen Zahlen und sei $A = \{ 2 \, x : x \in U \}$ die Menge der \emph{geraden} nat\"urlichen Zahlen.
% Die Frage, ob eine beliebige nat\"urliche Zahl $x \in U$ auch gerade ist (also ob $x \in A$), kann algorithmisch leicht in Zeit $\bigO(1)$ entschieden werden, indem wir $\ceil{\nicefrac{x}{2}} \cdot 2 \overset{?}{=} x$ auswerten und ggf. $1$ bzw. $0$ ausgeben.
% \end{example}

% \end{frame}

% \begin{frame}{Problemreduktion}
% Seien $A \subseteq U_1$ und $B \subseteq U_2$ und sei $A$ \emph{nicht} in $\bigO(g)$ entscheidbar.
% Wir wollen zeigen, dass $B$ \emph{ebenfalls nicht} in $\bigO(g)$ entschieden werden kann.

% Zwecks Erzeugung eines \alert{Widerspruchs} nehmen wir an, wir h\"atten einen Algorithmus $\mathbb{B}$, der $B$ in $\bigO(g)$ entscheidet.

% Wenn wir \emph{zeigen} können, dass wir jedes $a \in U_1$ mittels eines Algorithmus $\mathbb{T}$ in Zeit $\bigO(g)$ in ein $\mathbb{T}(a) \in U_2$ \alert{transformieren} können, sodass $a \in A \iff \mathbb{T}(a) \in B$, dann entscheidet folgender Algorithmus $A$ in Zeit $\bigO(g)$:
% \begin{align*}
%     &b \gets \mathbb{T}(a) \\
%     &\mathbf{return} \  \mathbb{B}(b)
% \end{align*}

% Das steht aber im Widerspruch dazu, dass $A$ \emph{nicht} in Zeit $\bigO(g)$ entschieden werden \emph{kann}.
% \end{frame}

\begin{frame}{Aufgaben zum Sortieren spezieller Arrays I}
\begin{task}
Sei $L$ eine Liste der L\"ange $n$ mit $L_i \leq L_{i + 2}$ f\"ur alle $1 \leq i \leq n - 2$.
Gibt es ein \emph{allgemeines Sortierverfahren}, das $L$ in Zeit $\bigO(n)$ sortiert?
\end{task}

\begin{solution}<2->
Betrachte die Teillisten $L' = \alglist{L_1, L_3, L_5, \dots}$ und $L'' = \alglist{L_2, L_4, L_6, \dots}$.
Diese k\"onnen wir in Zeit $\bigO(n)$ aus $L$ bestimmen und per Voraussetzung sind sie bereits sortiert.
Indem wir nun $L'$ und $L''$ in $\bigO(n)$ \emph{mergen}, haben wir $L$ in $\bigO(n)$ sortiert.
\end{solution}
\end{frame}

\begin{frame}{Aufgaben zum Sortieren spezieller Arrays II}
\begin{task}
Sei $L$ eine Liste der L\"ange $n$ mit $L_{\floor{\nicefrac{n}{2}} + 1} = L_{\floor{\nicefrac{n}{2}} + 2} = \dots = L_n$.
Gibt es ein \emph{allgemeines Sortierverfahren}, das $L$ in Zeit $\bigO(n)$ sortiert?
\end{task}

\begin{solution}<2->
Angenommen, es g\"abe einen Algorithmus $\mathbb{A}$, der \emph{solche} Listen in $\bigO(n)$ sortieren kann.
Dann k\"onnen wir in einer \alert{beliebigen} Liste $L$ der L\"ange $n$ einfach einen beliebigen Wert $a$ ans Ende anh\"angen -- und zwar $n$-mal -- und so in Zeit $\bigO(n)$ eine Liste $L'$ der L\"ange $2 \, n$ erhalten, die die Voraussetzungen von Algorithmus $\mathbb{A}$ erf\"ullt.

Wir sortieren nun $L'$ mit $\mathbb{A}$ in Zeit $\bigO(n)$ und l\"oschen anschlie{\ss}end die $n$ hinzugef\"ugten Vorkommen von $a$ wieder aus der sortierten Liste.

Somit haben wir $L$ in Zeit $\bigO(n)$ sortiert --- Widerspruch zu Theorem \ref{sort:thm:bound}!
\end{solution}
\end{frame}

\againframe<4>{general_sort}

\begin{frame}{Spezielle Sortierverfahren}
H\"aufig haben wir \alert{spezielle Kenntnisse} \"uber unsere Eingabedaten
\begin{itemize}
    \item Das Alter von Kunden in Jahren ist in $[0, 122]$
    \item Die K\"orpergr\"o{\ss}e von Patienten in Zentimetern ist in $[6, 272]$
    \item \emph{Hierarchische Strings}
    \begin{itemize}
        \item Raumnummern an der HU wie etwa \texttt{RUD 25, 4'112}
        \item Standorte in Bibliotheken wie etwa \texttt{ST 270 A149}
    \end{itemize}
    \item UNIX Nutzerrechte wie \texttt{666}
\end{itemize}
Wir k\"onnen hier nicht nur vergleichen, sondern auch die \alert{Struktur} der Datenpunkte ausnutzen.
\end{frame}

\begin{frame}{\"Uberblick \"uber spezielle Sortierverfahren}
In der Vorlesung wurden folgende \emph{spezielle} Sortierverfahren vorgestellt:
\begin{itemize}
    \item Radix Exchange Sort
    \item Bucket Sort
\end{itemize}
Im Folgenden f\"uhren wir \emph{Schreibtischtests} o.g. Verfahren durch.

\medskip

\textbf{Algorithmen folgen!}
\end{frame}

    \section{Suchverfahren}

\begin{frame}{Grundlagen}
\begin{block}{Problemstellung}
Gegeben eine \alert{sortierte} Liste $L = \alglist{a_1, \dots, a_n}$ \"uber einem Wertebereich $D$ und ein beliebiges Element $a \in D$.

Wir wollen die Frage beantworten, ob ein $i \in \{ 1, \dots, n \}$ mit $a_i = a$ existiert.
\end{block}
\begin{block}{Ansatz}<2->
Wir sprechen \"uber \alert{sortierte} Listen, k\"onnen also die \alert{Transitivit\"at} der $\leq_D$-Relation nutzen.

Das bedeutet, dass wir nach Vergleich eines Listenelements $a_i \neq a$ nur noch die Teilliste $L[1, i-1]$ bzw. $L[i+1, n]$ betrachten m\"ussen:
$$L = \alglist{\underbrace{a_1, \dots, a_{i-1}}_{\leq a_i}, a_i, \underbrace{a_{i+1}, \dots, a_n}_{\geq a_i}}$$
\end{block}
\end{frame}

\begin{frame}{Suchverfahren}
In der Vorlesung wurden folgende Suchverfahren vorgestellt:
\begin{itemize}
    \item Binary Search
    \item Fibonacci Search
    \item Interpolation Search
\end{itemize}
Die Verfahren unterscheiden sich im Wesentlichen durch das \alert{Kriterium}, nach dem das n\"achste zu vergleichende Element der Liste gew\"ahlt wird.
\begin{itemize}
    \item Ansonsten wird stets wie zuvor beschrieben die Transitivit\"at ausgenutzt
\end{itemize}

Im Folgenden führen wir \emph{Schreibtischtests} o.g. Verfahren durch.
\end{frame}

\begin{frame}{Binary Search}
\begin{columns}[T,onlytextwidth]
\begin{column}{0.42\textwidth}
\begin{algorithm}[H]
	\caption{Binary Search}
	\label{search:alg:binsearch}
	\DontPrintSemicolon
	\Input{Sortierte Liste $L$ \"uber $D$, \\ Element $a \in D$}
    \Output{Ist $a \in L$?}
    $\ell \gets 1, \quad r \gets n$\;
    \While{$\ell \leq r$}{
        $m \gets \ell + \floor{\frac{r-\ell}{2}}$\;
        \lIf{$a < L_m$}{$r \gets m - 1$}
        \lElseIf{$a > L_m$}{$\ell \gets m + 1$}
        \lElse{\Return{\True}}
    }
    \Return{\False}
\end{algorithm}
\end{column}
\begin{column}{0.53\textwidth}
\begin{task}<2->
Sei $L = \alglist{4, 4, 5, 17, 46, 58, 63}$.
Suchen Sie folgende Elemente in $L$ mittels des nebenstehenden Algorithmus:
\begin{tasks}
    \task $46$
    \task $3$
\end{tasks}
Geben Sie f\"ur jeden Durchlauf die Werte von $\ell$, $r$ und $m$ an!
\end{task}
\end{column}
\end{columns}
\end{frame}

\begin{frame}{Interpolation Search}
\begin{columns}[T,onlytextwidth]
\begin{column}{0.45\textwidth}
\begin{algorithm}[H]
	\caption{Interpolation Search}
	\label{search:alg:intsearch}
	\DontPrintSemicolon
	\Input{Sortierte Liste $L$ \"uber $D$, \\ Element $a \in D$}
    \Output{Ist $a \in L$?}
    $\ell \gets 1, \quad r \gets n$\;
    \While{$\ell \leq r$}{
        \only<1-2>{$m \gets \ell + \floor{(r - \ell) \cdot \red<2->{\frac{a - L_{\ell}}{L_r - L_{\ell}}}}$\;}
        \only<3->{
        \lIf{$L_r = L_{\ell}$}{$m \gets \ell$}
        \lElse{$m \gets \ell + \floor{(r - \ell) \cdot \frac{a - L_{\ell}}{L_r - L_{\ell}}}$}
        \BlankLine
        }
        \lIf{$a < L_m$}{$r \gets m - 1$}
        \lElseIf{$a > L_m$}{$\ell \gets m + 1$}
        \lElse{\Return{\True}}
    }
    \Return{\False}
\end{algorithm}
\end{column}
\begin{column}{0.5\textwidth}
\begin{task}<3->
Sei $L = \alglist{4, 4, 5, 17, 46, 58, 63}$.
Suchen Sie folgende Elemente in $L$ mittels des nebenstehenden Algorithmus:
\begin{tasks}
    \task $46$
    \task $3$
\end{tasks}
Geben Sie f\"ur jeden Durchlauf die Werte von $\ell$, $r$ und $m$ an!
\end{task}
\end{column}
\end{columns}
\end{frame}

\begin{frame}{Fibonacci-Zahlen}
\begin{definition}[Fibonacci-Zahlen]
Die Folge der \alert{Fibonacci-Zahlen} $(F_k)_{k \in \natpos}$ ist induktiv wie folgt definiert:
$$F_1 = 1, \quad F_2 = 1, \quad F_k = F_{k-2} + F_{k-1} \text{ f\"ur alle } k > 2.$$
\end{definition}

\begin{remark}
In den \"Ubungsaufgaben wird zudem gezeigt, dass $F_k \geq \nicefrac{1}{3} \cdot F_{k+2}$ f\"ur alle $k \geq 1$ gilt.
\end{remark}
\end{frame}

\begin{frame}{Fibonacci Search}
\scriptsize
\begin{algorithm}[H]
	\caption{Fibonacci Search}
	\label{search:alg:fibsearch}
	\DontPrintSemicolon
	\Input{Sortierte Liste $L$ \"uber $D$, Element $a \in D$}
    \Output{Ist $a \in L$?}
    $i \gets \min \{ k \in \natpos : F_k > n \}$\;
    $fib_3 \gets F_{i-3}, \quad fib_2 \gets F_{i-2}, \quad m \gets fib_2$\;
    \While{\True}{
        \If{$a < L_m$}{
            \lIf{$fib_2 = 1$}{\Return{\False}}
            \Else{
                $m \gets m - fib_3$\;
                $fib_2 \gets fib_2 - fib_3, \quad fib_3 \gets fib_3 - fib_2$\:
            }
        }
        \ElseIf{$a > L_m$}{
            \lIf{$fib_3 = 0$}{\Return{\False}}
            \Else{
                $m \gets m + fib_3$\;
                $t \gets fib_3, \quad fib_3 \gets fib_2 - fib_3, \quad fib_2 \gets t$\;
            }
        }
        \lElse{\Return{\True}}
    }
    \Return{\False}
\end{algorithm}
\end{frame}

    \section{Amortisierte Analyse}

\begin{frame}{Arrays}
\begin{definition}[Array]
Ein \alert{Array} $A = [ a_1, \dots, a_n ]$ der L\"ange $n$ repr\"asentiert eine Liste $L = \alglist{b_1, \dots, b_m}$ mit $m \leq n$, falls $a_i = b_i$ f\"ur $1 \leq i \leq m$ und $a_{m+1} = \dots = a_n = \square$.

Wir schreiben $\abs{A} \coloneqq \abs{L} = m$, um die Anzahl der \alert{belegten Slots} zu bezeichnen und $\norm{A} \coloneqq n$, um die \alert{Gesamtgr\"o{\ss}e} des Arrays zu bezeichnen.
Weiter sei $\nicefrac{\abs{A}}{\norm{A}}$ der \alert{F\"ullgrad}.
\end{definition}

\begin{example}
Sei $L = \alglist{\purple{8, 5, 12, 7, 9}}$, also $\abs{L} = 5$.
Das Array $A = [ \underbrace{\underbrace{\purple{8, 5, 12, 7, 9}}_{\abs{A} = 5}, \square, \square, \square}_{\norm{A} = 8} ]$ repr\"asentiert $L$.
\end{example}

\begin{remark}
Wir nehmen in der Regel an, dass $\norm{A} = 2^k$ f\"ur ein $k \in \nat$ ist.
\end{remark}
\end{frame}

\begin{frame}{Dynamische Arrays}
Ein \alert{dynamisches Array} repr\"asentiert eine Liste, an deren Ende "unbegrenzt" eingef\"ugt werden kann.
Gleichsam k\"onnen Elemente vom Ende wieder gel\"oscht werden.

Diese Operationen nennen wir im Folgenden \texttt{INS} f\"ur \alert{Einf\"ugung} und \texttt{DEL} f\"ur \alert{L\"oschung}.

Die Grundidee ist hierbei, \alert{nach Größenänderungen} stets einen \alert{F\"ullgrad} von $\geq \nicefrac{1}{2}$ sicherzustellen.
\begin{itemize}
    \item Das tun wir, indem wir bei Einf\"ugung des $(\norm{A}+1)$-ten Elements den verf\"ugbaren Speicherplatz \alert{verdoppeln}, 
    \item und bei L\"oschung des $\nicefrac{\norm{A}}{\red{4}}$-ten Elements den verf\"ugbaren Speicherplatz \alert{halbieren}.
\end{itemize}
\end{frame}

\begin{frame}{Operationen auf Dynamischen Arrays}
\begin{block}{Zu \texttt{INS}}
Sei $A = [ a_1, \dots, a_n ]$ ein Array, an dessen Ende wir ein neues Element $b$ \alert{einf\"ugen} wollen.
\begin{itemize}
    \item Falls $\abs{A} = \norm{A}$: Alloziere ein neues Array $A' = [ \underbrace{a_1, \dots, a_n}_{=A}, b, \underbrace{\square, \dots, \square}_{(n-1)\text{-mal}} ]$ mit \red{$\norm{A'} = 2\, n$}.
    \item Sonst f\"uge $b$ einfach in den n\"achsten freien Slot ein.
\end{itemize}
\end{block}

\begin{block}{Zu \texttt{DEL}}
Nun wollen wir das Element $a_m$ (mit $m = \abs{A}$) vom Ende von $A$ \alert{l\"oschen}.
\begin{itemize}
    \item Falls $\abs{A} = \red{\nicefrac{\norm{A}}{4}}$: Alloziere ein neues Array $A' = [ \underbrace{a_1, \dots, a_{m-1}}_{(\nicefrac{n}{4} - 1)\text{-mal}}, \underbrace{\square, \dots, \square}_{(\nicefrac{n}{4} + 1)\text{-mal}} ]$ mit \red{$\norm{A'} = \nicefrac{n}{2}$}.
    \item Sonst l\"osche einfach $a_m$.
\end{itemize}
\end{block}
\end{frame}

\begin{frame}{Analyse der Operationen auf Dynamischen Arrays}
\begin{block}{Analyse}
Eine \texttt{INS}-Operation kostet im Worst-Case $\bigO(n)$, da s\"amtliche $n-1$ Elemente verschoben und erst dann die Einf\"ugung get\"atigt werden kann.

Analog gilt f\"ur die \texttt{DEL}-Operation, dass das Verschieben von $\nicefrac{n}{4}-1$ Elementen in $\bigO(n)$ ist.
\end{block}

\begin{remark}
Folglich kostet eine Sequenz $n$ direkt aufeinanderfolgender \texttt{INS}-Operationen $\bigO(n^2)$.
\end{remark}

\begin{block}{Problem}
Wir betrachten die Worst-Case Lauzeit \alert{jeder einzelnen Operation} --- diese Sichtweise ist hier \alert{zu pessimistisch}.
\end{block}

\begin{block}{L\"osung}
Betrachte stattdessen \alert{Sequenzen von Operationen} und deren Gesamtlaufzeit.
\end{block}
\end{frame}

\begin{frame}{Dynamische Arrays --- Beispiel}
\begin{example}
Betrachte die Sequenz von Operationen $o_1, \dots, o_n$ mit $o_1 = \dots = o_n = \text{\texttt{INS}}$.

\begin{columns}[T]
\begin{column}{0.8\textwidth}
\setlength{\leftskip}{1em}
\only<2->{
In nebenstehender "Verteilung" der Laufzeiten f\"allt auf, dass \emphred{$o_i$} genau dann \emphred{nicht konstant} ist, wenn \emphorange<2>{$i-1$} eine \emphorange<2>{Zweierpotenz} ist.
}

\medskip

\only<3->{
Seien $c_1, \dots, c_n$ die \alert{tats\"achlichen Kosten} der $i$-ten Operation, wobei 
$$c_i = \begin{cases} \red{i} & \red{\text{falls } i-1 \text{ eine Zweierpotenz ist}} \\ \teal{1} & \text{\teal{sonst}.}  \end{cases}$$
}

\only<4->{
Dann belaufen sich die \alert{Gesamtkosten} aller $n$ Operationen auf
$${\displaystyle
\sum_{i=1}^n c_i 
\leq \teal{n} + \red{\sum_{i=0}^{\floor{\log n}} 2^i} 
\leq \teal{n} + \red{\sum_{i=0}^{\floor{\log n}} \frac{n}{2^i}}
\leq \teal{n} + \red{2 \, n} = 3 \, n
}.$$
}
\end{column}
\begin{column}<2->{0.2\textwidth}
\centering
\begin{tabular}{cc}
    $\mathbf{i}$ & \textbf{Laufzeit} \\
    \orange<2>{\teal<3->{1}} & \teal<3->{$\bigO(1)$} \\
    \red<2->{2} & \red<2->{$\bigO(i)$} \\
    \red<2->{3} & \red<2->{$\bigO(i)$} \\
    \orange<2>{\teal<3->{4}} & \teal<3->{$\bigO(1)$} \\
    \red<2->{5} & \red<2->{$\bigO(i)$} \\
    \teal<3->{6} & \teal<3->{$\bigO(1)$} \\
    \teal<3->{7} & \teal<3->{$\bigO(1)$} \\
    \orange<2>{\teal<3->{8}} & \teal<3->{$\bigO(1)$} \\
    \red<2->{9} & \red<2->{$\bigO(i)$} \\
    $\vdots$ & $\vdots$
\end{tabular}
\end{column}
\end{columns}
\end{example}
\end{frame}

\begin{frame}{Dynamische Arrays --- Zwischenstand}
\begin{block}{Ergebnis}
Wir haben gezeigt, dass jede Sequenz von $n$ Einf\"ugungen insgesamt $\leq 3 \, n$ Instruktionen ben\"otigt --- jede einzelne Einf\"ugung ben\"otigt also \alert{amortisiert} $3$ Instruktionen.
\end{block}

\begin{block}{Intuition}
Die amortisierten Kosten je Einf\"ugung in H\"ohe von $3$ kann man wie folgt veranschaulichen:
\begin{itemize}
    \item Wir nutzen eine Einheit, um \alert{das neue Element selbst einzuf\"ugen}
    \item Eine weitere Einheit nutzen wir, um f\"ur \alert{dessen sp\"atere Verschiebung} zu zahlen
    \item Schlie{\ss}lich ben\"otigen wir eine Einheit, um f\"ur's Verschieben von Elementen zu zahlen, die sich bereits \alert{vor der letzten Vergr\"o{\ss}erung} in der Liste befanden
\end{itemize}
\end{block}
\end{frame}

\begin{frame}{\"Uberblick}
Zur Analyse der Gesamtlaufzeit einer \alert{Sequenz von Operationen} stehen folgende Verfahren zur Verf\"ugung:
\begin{itemize}
    \item Aggregationsmethode
    \item Account-Methode
    \item Potentialmethode
\end{itemize}
Da diese paarweise \alert{\"aquivalent} zueinander sind, k\"onnen wir schlicht das zum jeweiligen Problem passende Verfahren w\"ahlen.

\begin{quote}
    Der Grundgedanke ist stets, dass sich der zus\"atzliche Aufwand der \alert{teureren}, aber daf\"ur \alert{selteneren} Operationen auf die \"ubrigen Operationen \alert{verteilt}.
\end{quote}

Im Folgenden betrachten wir stellvertretend die \alert{Potentialmethode}.
\end{frame}

\begin{frame}{Die Potentialmethode}
\begin{definition}[Potentialmethode]
Sei $D$ eine Datenstruktur und $o_1, \dots, o_m$ eine Sequenz von Operationen auf $D$.
Weiterhin sei $D_0 \coloneqq D$ und $D_i \coloneqq o_i(D_{i-1})$ f\"ur $1 \leq i \leq m$ der \alert{Zustand} von $D$ \emph{nach} Anwendung der $i$-ten Operation $o_i$.

\begin{itemize}
    \item Eine \alert{Potentialfunktion} $\Phi$ weist jedem Zustand $D_i$ ein \alert{Potential} $\Phi(D_i) \in \realpos$ zu.
    \item Die \alert{amortisierten Kosten} der $i$-ten Operation sind $\hat{c_i} \coloneqq c_i + \Phi(D_i) - \Phi(D_{i-1})$, wobei $c_i$ die \alert{tats\"achlichen Kosten} sind.
\end{itemize}
\end{definition}

\begin{remark}
Der Einfachheit halber werden wir meist annehmen, dass $\Phi(D_0) = 0$ ist.
\end{remark}

\begin{theorem}\label{arm:thm:equiv}
\vspace*{-6pt}
Falls eine Potentialfunktion $\Phi$ mit $\Phi(D_n) \geq \Phi(D_0)$ existiert, so ist $\sum_{i=1}^n \hat{c_i} \geq \sum_{i=1}^n c_i$.
\end{theorem}
\end{frame}

\begin{frame}{Amortisierte vs. Tats\"achliche Kosten}
\begin{proof}[Beweis von Theorem \ref{arm:thm:equiv}]
Es ist leicht zu zeigen, dass die Summe der amortisieren Kosten eine \alert{\red<2->{obere Schranke} f\"ur die tats\"achlichen Kosten} der Sequenz darstellt:
\begin{align*}
    \red<2->{\sum_{i=1}^n \hat{c_i}} &= \sum_{i=1}^n (c_i + \Phi(D_i) - \Phi(D_{i-1})) \\
    &= \sum_{i=1}^n c_i + \underbrace{\sum_{i=1}^n (\Phi(D_i) - \Phi(D_{i-1}))}_{\text{Teleskopsumme}} \\
    &= \sum_{i=1}^n c_i + \underbrace{\Phi(D_n) - \Phi(D_0)}_{\geq 0} \red<2->{\geq \sum_{i=1}^n c_i}
\end{align*}
\end{proof}
\end{frame}

\begin{frame}{Amortisierte Analyse Dynamischer Arrays \, I}
\begin{theorem}
Jede Sequenz von $n$ \texttt{INS}-Operationen auf einem dynamischen Array $A$ ist in $\bigO(n)$.
\end{theorem}
\begin{block}{Beweis.}
Wir nutzen zum Beweis die Potentialmethode.
Sei hierf\"ur $\Phi(A_i) \coloneqq 2 \cdot \abs{A} - \norm{A}$.
Insbesondere gilt somit f\"ur ein \emph{volles} Array (d.h. $\abs{A} = \norm{A}$), dass $\Phi(A_i) = 2 \cdot \norm{A} - \norm{A} = \norm{A}$.

F\"ur die Ermittlung der amortisierten Kosten unterscheiden wir zwei F\"alle.
\begin{itemize}
    \item Das Array enth\"alt noch einen freien Slot
    \item Das Array ist voll und wir m\"ussen zun\"achst ein neues Array allozieren
\end{itemize}
\end{block}
\end{frame}

\begin{frame}{Amortisierte Analyse Dynamischer Arrays \, II}
\begin{block}{Fall 1: $\abs{A_{i-1}} < \norm{A_{i-1}}$}
Das Array enth\"alt also noch einen freien Slot f\"ur das neue Element, sodass sich die Gr\"o{\ss}e von $A$ \emph{nicht} ver\"andert, d.h. $\norm{A_i} = \norm{A_{i-1}}$.

Einzig die Anzahl der belegten Slots erh\"oht sich um eins, d.h. $\abs{A_i} = \abs{A_{i-1}} + 1$.
Folglich belaufen sich die amortisierten Kosten der $i$-ten Operation f\"ur diesen Fall auf:
\begin{align*}
    \hat{c_i} &= c_i + \Phi(A_i) - \Phi(A_{i-1}) \\
    &= c_i + (2 \cdot \underbrace{\abs{A_i}}_{= \abs{A_{i-1}} + 1} - \underbrace{\norm{A_i}}_{=\norm{A_{i-1}}}) - (2 \cdot \abs{A_{i-1}} - \norm{A_{i-1}}) \\
    &= c_i + 2 \cdot \abs{A_{i-1}} + 2 - \norm{A_{i-1}} - 2 \cdot \abs{A_{i-1}} + \norm{A_{i-1}} \\
    &= \underbrace{c_i}_{=1} + 2 \\
    &= 3
\end{align*}
\end{block}
\end{frame}

\begin{frame}{Amortisierte Analyse Dynamischer Arrays \, III}
\begin{block}{Fall 2: $\abs{A_{i-1}} = \norm{A_{i-1}}$}
Wir m\"ussen also zun\"achst ein neues Array der Gr\"o{\ss}e $2 \cdot \norm{A_{i-1}}$ allozieren und die bisherigen $\norm{A_{i-1}}$ Elemente verschieben.

Die Anzahl der belegten Slots erh\"oht sich wiederum um eins, d.h. $\abs{A_i} = \abs{A_{i-1}} + 1$.
Insgesamt belaufen sich die amortisierten Kosten der $i$-ten Operation f\"ur diesen Fall auf:
\begin{align*}
    \hat{c_i} &= c_i + \Phi(A_i) - \Phi(A_{i-1}) \\
    &= c_i + (2 \cdot \underbrace{\abs{A_i}}_{= \abs{A_{i-1}} + 1} - \underbrace{\norm{A_i}}_{=2\cdot \norm{A_{i-1}}}) - (2 \cdot \abs{A_{i-1}} - \norm{A_{i-1}}) \\
    &= c_i + 2 \cdot \abs{A_{i-1}} + 2 - 2 \cdot \norm{A_{i-1}} - 2 \cdot \abs{A_{i-1}} + \norm{A_{i-1}} \\
    &= \underbrace{c_i}_{=\norm{A_{i-1}} + 1} + 2 - \norm{A_{i-1}} \\
    &= 3
\end{align*}
\end{block}
\end{frame}

\begin{frame}{Amortisierte Analyse Dynamischer Arrays \, IV}
Die \emph{tats\"achlichen} Kosten f\"ur alle $n$ Operationen liegen gem\"a{\ss} Theorem \ref{arm:thm:equiv} also bei $\leq 3 \, n \in \bigO(n)$.
Dies beendet den Beweis. \qed

\begin{remark}
Auf \"ahnliche Weise kann gezeigt werden, dass die Aussage auch dann noch gilt, wenn man auch \texttt{DEL}-Operationen zul\"asst.
Details dazu finden sich bspw. in Kapitel 17.4.2 von

\begin{quote}
    Cormen, Leiserson, Rivest, Stein: Introduction to Algorithms, 2nd Edition, MIT Press
\end{quote}
\end{remark}
\end{frame}

    \section{Heaps}

\begin{frame}{Priorit\"atswarteschlangen}
\begin{block}<+->{Grundidee}
Wir wollen schnell auf das aktuell h\"ochstpriorisierte Element einer \alert{ver\"anderlichen} Liste zugreifen k\"onnen.

Genauer gesagt soll die gesuchte Datenstruktur folgende Operationen unterst\"utzen:
\begin{itemize}
    \item \alert{Finden} und \alert{Entfernen} des aktuell h\"ochstpriorisierten Elements
    \item \alert{Einf\"ugen} neuer Elemente
\end{itemize}
\end{block}

\begin{block}<+->{Problem}
F\"ur sortierte/unsortierte Arrays/Linked-Lists ist stets mindestens eine der gew\"unschten Operationen in $\red{\bigO(n)}$.
\end{block}

\begin{block}<+->{L\"osung}
Wir organisieren unsere Daten in einem \alert{Baum}, den wir durch ein \alert{Array} repr\"asentieren.
\end{block}
\end{frame}

\begin{frame}{Heaps}
\begin{definition}\label{heaps:def:heap}\vspace*{-6pt}
Ein \alert{Heap} ist ein \emph{beinahe vollst\"andiger Bin\"arbaum} $T = (V, E, r)$ mit Beschriftungsfunktion $\pi: V \to \nat$, f\"ur den gilt:
\begin{quote}
    Ist $\purple<2->{v} \in V$ ein innerer Knoten mit Kindern $\teal<2->{v_{\ell}}, \red<2->{v_r}$, so ist $\pi(\purple<2->{v}) \leq \pi(\teal<2->{v_{\ell}})$ und $\pi(\purple<2->{v}) \leq \pi(\red<2->{v_r})$.
\end{quote}
\end{definition}

\begin{block}<2->{Skizze zu Definition \ref{heaps:def:heap}}
\begin{tikzpicture}[every edge/.append style={draw, -}]
    \node[amaranth, fill=amaranth!7] (A) at (0,0) {4};
    \node[tealblue, fill=tealblue!7] (B) at (-1,-1) {7};
    \node[madderlake, fill=madderlake!7] (C) at (1,-1) {5};
    
    \node[draw=none] (D) at (-0.6,0.6) {};
    \node[draw=none] (E) at (-1.6,-1.6) {};
    \node[draw=none] (F) at (-0.4,-1.6) {};
    \node[draw=none] (G) at (0.4,-1.6) {};
    \node[draw=none] (H) at (1.6,-1.6) {};
    
    \path (A) edge (B);
    \path (A) edge (C);
    \path (A) edge (D);
    \path (B) edge (E);
    \path (B) edge (F);
    \path (C) edge (G);
    \path (C) edge (H);
\end{tikzpicture}
\end{block}
\end{frame}

\begin{frame}{Operationen auf Heaps}
Im Folgenden betrachten wir Operationen zum Erstellen und Verwalten von Heaps:
\begin{itemize}
    \item \lightblue{\textsc{BottomUpSiftDown}} zum Aufbau einer Heap-Struktur
    \item \lightblue{\textsc{Add}} zum Einf\"ugen eines neuen Elements
    \item \lightblue{\textsc{ExtractMin}} zum Entfernen des minimalen Elements
\end{itemize}

\begin{remark}<2->
S\"amtliche o.g. Operationen lassen sich durch Anpassung einiger Vergleiche leicht auf sog. Max-Heaps \"ubertragen.
\end{remark}
\end{frame}

% \begin{frame}{Prequel: Die Heapify Operation}
% \begin{columns}[T]
% \begin{column}{0.4\textwidth}
% \begin{algorithm}[H]
% 	\caption{Heapify}
% 	\label{heap:alg:extractmin}
% 	\DontPrintSemicolon
% 	\Input{Dynamisches Array $A$}
%     \Output{Minimales Element aus $A$}
% 	$min \gets A_1$\;
% 	Vertausche $A_1$ und $A_n$\;
% 	L\"osche $A_n$\;
% 	\textsc{Heapify($A$)}\;
%     \Return{$min$}\;
% \end{algorithm}
% \end{column}
% \begin{column}{0.55\textwidth}

% \end{column}
% \end{columns}
% \end{frame}

% \begin{frame}{Entfernen des minimalen Elements}
% \begin{columns}[T]
% \begin{column}{0.4\textwidth}
% \begin{algorithm}[H]
% 	\caption{ExtractMin}
% 	\label{heap:alg:extractmin}
% 	\DontPrintSemicolon
% 	\Input{Dynamisches Array $A$}
%     \Output{Minimales Element aus $A$}
% 	$min \gets A_1$\;
% 	Vertausche $A_1$ und $A_n$\;
% 	L\"osche $A_n$\;
% 	\textsc{Heapify($A$)}\;
%     \Return{$min$}\;
% \end{algorithm}
% \end{column}
% \begin{column}{0.55\textwidth}
% Wir vertauschen also die Wurzel mit dem letzten Blatt, stellen anschlie{\ss}end die Heap-Eigenschaft wieder her und geben den Wert der vormaligen Wurzel zur\"uck.

% \medskip

% Die Laufzeit liegt somit offensichtlich \alert{amortisiert} in $\bigO(\log n)$.
% \begin{itemize}
%     \item Amortisiert deshalb, weil das L\"oschen des letzten Elements aus einem \emph{dynamischen} Array im Worst-Case $\bigO(n)$ kostet
% \end{itemize}
% \end{column}
% \end{columns}
% \end{frame}

% \begin{frame}{Einf\"ugen eines neuen Elements}
% \begin{columns}[T]
% \begin{column}{0.4\textwidth}
% \begin{algorithm}[H]
% 	\caption{Add}
% 	\label{heap:alg:add}
% 	\DontPrintSemicolon
% 	\Input{Dynamisches Array $A$,\\Neuer Wert $a$}
%     \Output{Dynamisches Array $A'$}
%     $A_{n+1} \gets a$\;
%     $c \gets n + 1, \quad p \gets \floor{\nicefrac{c}{2}}$\;
%     \While{$p > 0$ \And $A_c < A_p$}{
%         Vertausche $A_c$ und $A_p$\;
%         $c \gets p, \quad p \gets \floor{\nicefrac{c}{2}}$\;
%     }
% \end{algorithm}
% \end{column}
% \begin{column}{0.55\textwidth}
% Wir machen also praktisch "BubbleSort" auf dem \emph{eindeutigen} Pfad von dem neu eingef\"ugten Blatt mit Wert $a$ und der Wurzel.

% \medskip

% Da heapgeordnete Arrays beinahe vollst\"andig sind, ist die L\"ange dieses Pfades in $\bigO(\log n)$ und somit auch die Laufzeit des Algorithmus'.
% \end{column}
% \end{columns}
% \end{frame}

% \begin{frame}{Aufbau einer Heap-Struktur}
    
% \end{frame}
    \section{Graphen \& B\"aume}

\begin{frame}<1-4>[label=graph]
\frametitle<1-4>{Graphen}
\frametitle<5>{Zur Erinnerung: Graphen}
\begin{columns}[T]
\begin{column}{0.6\textwidth}
\begin{definition}[Graph]
\only<+->{
Ein \alert{(ungerichteter) Graph} ist ein Paar $G = (\red<1>{V}, \orange<1>{E})$ mit einer Menge von \alert{\emphred<1>{Knoten}} $\red<1>{V}$ und einer Menge von \alert{\emphorange<1>{Kanten}} $\orange<1>{E} \subseteq \{ \{v, w \} \subseteq \red<1>{V} : v \neq w \}$.
}

\only<+->{
Ein \alert{\emphred<2>{Pfad}} der L\"ange $k$ ist eine Folge \emph{paarweise verschiedener} Knoten $\red<2>{v_0, \dots, v_k} \in V$ mit $\orange<2>{\{ v_i, v_{i+1} \}} \in E$ f\"ur $0 \leq i < k$.
}

\only<+->{
Ist $v_1, \dots, v_k$ ein Pfad und ist $\purple<3>{\{ v_k, v_1 \}} \in E$, so nennen wir die Folge \emphred<3>{$v_1, \dots, v_k, v_1$} einen \alert{\emphred<3>{Kreis}} der L\"ange $k$.
}

\only<+->{
Ein Graph ist \alert{zusammenh\"angend}, falls f\"ur alle $(v, w) \in V \times V$ gilt, dass es einen Pfad $v_1, \dots, v_k$ mit $v_1 = v$ und $v_k = w$ gibt.
}
\end{definition}
\end{column}
\begin{column}{0.35\textwidth}
\vspace{0.8cm}
\begin{tikzpicture}[every edge/.append style={draw, -}]
    \tikzstyle{highlightred}=[madderlake, fill=madderlake!7]
    \tikzstyle{highlightorange}=[mandarin, fill=mandarin!7]

    \node[onslide={<1-3>{highlightred}}] (A) at (0,0) {A};
    \node[onslide={<1-3>{highlightred}}] (B) at (0.5,2) {B};
    \node[onslide={<1>{highlightred}}] (C) at (2,2.5) {C};
    \node[onslide={<1-3>{highlightred}}] (D) at (1.7,-0.2) {D};
    \node[onslide={<1>{highlightred}}] (E) at (2,-2) {E};
    \node[onslide={<1-3>{highlightred}}] (F) at (3,1.3) {F} ;
    
    \path[onslide={<1-3>{highlightorange}}] (A) edge (B);
    \path[onslide={<1>{highlightorange}}] (B) edge (C);
    \path[onslide={<1-3>{highlightorange}}] (A) edge (D);
    \path[onslide={<1>{highlightorange}}] (D) edge (C);
    \path[onslide={<1>{highlightorange}}] (A) edge (E);
    \path[onslide={<1>{highlightorange}}] (D) edge (E);
    \path[onslide={<1-3>{highlightorange}}] (D) edge (F);
    \path[onslide={<1>{highlightorange}}] (C) edge (F);
    \path[onslide={<1>{highlightorange}}] (C) edge (A); 
    \path[onslide={<1>{highlightorange}}] (E) edge (F);
    \only<3>{\path[amaranth, dashed] (F) edge (B);}
\end{tikzpicture}
\end{column}
\end{columns}
\pause
\end{frame}

\begin{frame}{Repr\"asentation von Graphen}
\begin{definition}
Sei $G = (V, E)$ ein Graph.
F\"ur alle Knoten $v \in V$ nennen wir $N(v) \coloneqq \{ w \in V : \{v, w\} \in E \}$ die \alert{Nachbarschaft} von $v$.
\end{definition}

\begin{block}{Adjazenzmatrix}
Wir k\"onnen einen Graphen $G = (V, E)$ mit $n \coloneqq \abs{V}$ und o.B.d.A. mit $V = \{ 1, \dots, n \}$ durch eine \alert{Adjazenzmatrix} $M_G = (m_{i,j})_{1 \leq i < j \leq n}$ mit $m_{i, j} = 1 \iff \{ i, j \} \in E$ darstellen.
\end{block}

\begin{block}{Adjazenzlisten}
Alternativ k\"onnen wir $G$ auch durch Folge von \alert{Adjazenzlisten} $N_1, \dots, N_n$ mit $N_i = N(i)$ f\"ur $1 \leq i \leq n$ darstellen, d.h. die $i$-te Liste enth\"alt genau die Nachbarn von $i$.
\end{block}
\end{frame}

\begin{frame}{Vergleich von Adjazenzmatrix und Adjazenzlisten}
Der Vorteil von Adjazenzmatrizen gegen\"uber Adjazenzlisten ist, dass f\"ur je zwei Knoten $v, w \in V$ in Zeit $\bigO(1)$ \"uberpr\"uft werden kann, ob $\{ v, w \} \in E$ ist --- f\"ur Adjazenzlisten ist dieser Test in $\bigO(n)$ bzw. $\bigO(\log n)$ falls die Listen sortiert sind.

Andererseits haben Adjazenzlisten den Vorteil, die Nachbarschaft eines Knotens $v$ in Zeit $\bigO(\abs{N(v)})$ iterieren zu k\"onnen, w\"ahrend dies f\"ur Adjazenzmatrizen in $\bigO(n)$ ist.

Allgemein gilt, dass sich die Vorteile der Adjazenzliste verlieren, je \alert{dichter} der Graph ist, d.h. je mehr Kanten er im Verh\"altnis zu seiner Knotenanzahl hat.
\begin{itemize}
    \item Beispielsweise gilt f\"ur den \alert{vollst\"andigen Graphen} $K_n = (V, E)$ mit $\abs{V} = n$ und $E = \{ \{ v, w \} : v, w \in V \}$, d.h. $m = \frac{n \cdot (n-1)}{2}$, dass $\abs{N(v)} = n - 1$, sodass auch die Adjazenzmatrix eine Iteration der Nachbarn in Zeit $\bigO(\abs{N(v)})$ erlaubt.
\end{itemize}
\end{frame}

\begin{frame}{Breitensuche \, I}
\begin{columns}[T]
\begin{column}<1->{0.45\textwidth}
\vspace*{-5pt}
\begin{algorithm}[H]
	\caption{Breitensuche (BFS)}
	\label{trees:alg:bfs}
	\vspace*{-6pt}
	\SetKwFunction{Enqueue}{enqueue}
	\SetKwFunction{Dequeue}{dequeue}
	\SetKwFunction{Visit}{visit}
	\DontPrintSemicolon
	\Input{Graph $G = (V, E)$, Pivot $p \in V$}
    $visited \gets \emptyset$\;
    $frontier \gets \{ p \}$\;
    \While{$frontier \neq \emptyset$}{
        $v \gets $ \Dequeue{$frontier$}\;
        \If{$v \not \in visited$}{
            \Visit{$v$}\;
            $visited \gets visited \cup \{ v \}$\;
            \ForEach{$w \in N(v)$}{
                \Enqueue{$frontier$, $w$}\;
            }
        }
    }
\end{algorithm}
\end{column}
\begin{column}<2->{0.5\textwidth}
Breitensuche ist eine grundlegende Strategie zum systematischen Durchlaufen eines Graphen.

\medskip

In \emph{visited} merken wir uns bereits besuchte Knoten, w\"ahrend \emph{frontier} eine \alert{Queue} \emph{eventuell} noch zu besuchender Knoten ist.

\medskip

Sobald wir einen Knoten $v$ aus \emph{frontier} entfernen, pr\"ufen wir zun\"achst, ob er bereits besucht wurde.
Falls nicht, f\"uhren wir eine nutzerdefinierte Funktion \texttt{visit} (bspw. \texttt{print}) aus und markieren $v$ als besucht.
Abschlie{\ss}end f\"ugen wir alle Nachbarn von $v$ zu \emph{frontier} hinzu.
\end{column}
\end{columns}
\end{frame}

\begin{frame}{Breitensuche \, II}
\begin{columns}[T]
\begin{column}{0.45\textwidth}
\vspace*{-5pt}
\addtocounter{algocf}{-1}
\begin{algorithm}[H]
	\caption{Breitensuche (BFS)}
	\SetKwFunction{Enqueue}{enqueue}
	\SetKwFunction{Dequeue}{dequeue}
	\SetKwFunction{Visit}{visit}
	\DontPrintSemicolon
	\Input{Graph $G = (V, E)$, Pivot $p \in V$}
    $visited \gets \emptyset$\;
    $frontier \gets \{ p \}$\;
    \While{$frontier \neq \emptyset$}{
        $v \gets $ \Dequeue{$frontier$}\;
        \If{$v \not \in visited$}{
            \Visit{$v$}\;
            $visited \gets visited \cup \{ v \}$\;
            \ForEach{$w \in N(v)$}{
                \Enqueue{$frontier$, $w$}\;
            }
        }
    }
\end{algorithm}
\end{column}
\begin{column}{0.5\textwidth}
Hinsichtlich der Laufzeit betrachten wir hier nur \emph{zusammenh\"angende} Graphen $G = (V, E)$ mit $n \coloneqq \abs{V}$ und $m \coloneqq \abs{E}$.
F\"ur diese gilt $n \leq m + 1$.

\medskip

Weiterhin nehmen wir an, dass \texttt{visit} in $\bigO(1)$ liegt und dass $G$ durch Adjazenzlisten repr\"asentiert ist, sodass $N(v)$ in $\bigO(1)$ ist.

\medskip

Dann ist klar, dass \algref{trees:alg:bfs} in Zeit $\bigO(m)$ ist, da alle $n < m$ Knoten und $m$ Kanten h\"ochstens zwei mal betrachtet werden (hier ohne Beweis).
\end{column}
\end{columns}
\end{frame}

\begin{frame}{Tiefensuche}
\begin{columns}[T]
\begin{column}<1->{0.45\textwidth}
\vspace*{-5pt}
\begin{algorithm}[H]
	\caption{Tiefensuche (DFS)}
	\label{trees:alg:dfs}
	\vspace*{-6pt}
	\SetKwFunction{Push}{\emphred<2>{push}}
	\SetKwFunction{Pop}{\emphred<2>{pop}}
	\SetKwFunction{Visit}{visit}
	\DontPrintSemicolon
	\Input{Graph $G = (V, E)$, Pivot $p \in V$}
    $visited \gets \emptyset$\;
    $frontier \gets \{ p \}$\;
    \While{$frontier \neq \emptyset$}{
        $v \gets $ \Pop{$frontier$}\;
        \If{$v \not \in visited$}{
            \Visit{$v$}\;
            $visited \gets visited \cup \{ v \}$\;
            \ForEach{$w \in N(v)$}{
                \Push{$frontier$, $w$}\;
            }
        }
    }
\end{algorithm}
\end{column}
\begin{column}<2->{0.5\textwidth}
Einziger Unterschied zu \algref{trees:alg:bfs} ist die Verwendung eines \emphred{Stacks} anstelle einer Queue.

\medskip

Entsprechend behalten auch die Aussagen bzgl. der Laufzeit ihre G\"ultigkeit.
\end{column}
\end{columns}
\end{frame}

\begin{frame}<1-3>[label=diameter]
\frametitle<1-3>{Der Durchmesser eines Graphen}
\frametitle<4>{Zur Erinnerung: Der Durchmesser eines Graphen}
\begin{columns}[T]
\begin{column}{0.6\textwidth}
\begin{definition}<+->[Durchmesser]
Der \emphorange<1>{Durchmesser} eines Graphen $G$ entspricht der L\"ange des \emph{l\"angsten} \emphorange<1>{Pfades} in $G$.
\end{definition}

\medskip

\only<+->{
\begin{quote}
    Dieser Parameter ist beispielsweise bei der Analyse von Netzwerken von gro{\ss}er Bedeutung.
\end{quote}
}

\begin{theorem}<+->
Im Allgemeinen liegt die Laufzeit der Berechnung des Durchmessers einen Graphen $G = (V, E)$ mit $n \coloneqq \abs{V}$ in $\bigO(n^3)$.
\end{theorem}
\end{column}
\begin{column}{0.35\textwidth}
\vspace{0.8cm}
\begin{tikzpicture}[every edge/.append style={draw, -}]
    \tikzstyle{highlightred}=[madderlake, fill=madderlake!7]
    \tikzstyle{highlightorange}=[mandarin, fill=mandarin!7]

    \node[onslide={<1>{highlightred}}] (A) at (0,0) {A};
    \node[onslide={<1>{highlightred}}] (B) at (0.5,2) {B};
    \node[onslide={<1>{highlightred}}] (C) at (2,2.5) {C};
    \node[onslide={<1>{highlightred}}] (D) at (1.7,-0.2) {D};
    \node[onslide={<1>{highlightred}}] (E) at (2,-2) {E};
    \node[onslide={<1>{highlightred}}] (F) at (3,1.3) {F} ;
    
    \path[onslide={<1>{highlightorange}}] (A) edge (B);
    \path[onslide={<1>{highlightorange}}] (B) edge (C);
    \path[] (A) edge (D);
    \path[onslide={<1>{highlightorange}}] (D) edge (C);
    \path[] (A) edge (E);
    \path[onslide={<1>{highlightorange}}] (D) edge (E);
    \path[] (D) edge (F);
    \path[] (C) edge (F);
    \path[] (C) edge (A); 
    \path[onslide={<1>{highlightorange}}] (E) edge (F);
\end{tikzpicture}
\end{column}
\end{columns}
\pause
\end{frame}

\againframe<5>{graph}

\begin{frame}{B\"aume}
\begin{columns}[T]
\begin{column}{0.6\textwidth}
\begin{definition}[Baum]
\only<1->{Ein \alert{Baum} ist ein \emph{zusammenh\"angender}, \emph{kreisfreier} Graph.}
\only<2->{Ein \alert{gewurzelter Baum} ist ein Tripel $T = (\red<2>{V}, \orange<2>{E}, \purple<2>{r})$ mit \alert{Wurzel} $\purple<2>{r} \in \red<2>{V}$, wobei $(\red<2>{V}, \orange<2>{E})$ ein Baum ist.}

\only<3->{Die \alert{Tiefe} $d(\red<3>{v})$ eines Knotens $\red<3>{v} \in V$ in $T$ ist die L\"ange des (eindeutigen) \emphorange<3>{Pfades} von $\red<3>{v}$ zur Wurzel $\purple<3>{r}$.}

\only<4->{In einer Kante $\orange<4>{\{ v, w \}} \in E$ mit $d(\teal<4>{w}) = d(\red<4>{v}) + 1$ nennen wir $\teal<4>{w}$ ein \alert{Kind} von $\red<4>{v}$ bzw. $\red<4>{v}$ \emph{den} \alert{Elter} von $\teal<4>{w}$.}

\only<5->{Knoten ohne Kinder sind \alert{\emphteal<5>{Bl\"atter}}, alle anderen sind \alert{\emphred<5>{innere Knoten}}.}
\end{definition}
\end{column}
\begin{column}{0.35\textwidth}
\vspace{0.8cm}
\begin{tikzpicture}[every edge/.append style={draw, -}]
    \tikzstyle{highlightred}=[madderlake, fill=madderlake!7]
    \tikzstyle{highlightteal}=[tealblue, fill=tealblue!7]
    \tikzstyle{highlightorange}=[mandarin]
    \tikzstyle{highlightpurple}=[amaranth, fill=amaranth!7]

    \node[onslide={<2,3>{highlightpurple}}, onslide={<5>{highlightred}}] (A) at (0,0) {A};
    \node[onslide={<2>{highlightred}}, onslide={<5>{highlightteal}}] (B) at (0.5,2) {B};
    \node[onslide={<2, 4, 5>{highlightred}}] (C) at (2,2.5) {C};
    \node[onslide={<2, 5>{highlightred}}] (D) at (1.7,-0.2) {D};
    \node[onslide={<2, 3, 5>{highlightred}}] (E) at (2,-2) {E};
    \node[onslide={<2>{highlightred}}, onslide={<4, 5>{highlightteal}}] (F) at (3,1.3) {F};
    \node[onslide={<2>{highlightred}}, onslide={<5>{highlightteal}}] (G) at (3.2,-1.2) {G};
    
    \path[onslide={<2>{highlightorange}}] (A) edge (B);
    \path[onslide={<2, 3>{highlightorange}}] (A) edge (D);
    \path[onslide={<2, 3>{highlightorange}}] (D) edge (E);
    \path[onslide={<2, 4>{highlightorange}}] (C) edge (F);
    \path[onslide={<2>{highlightorange}}] (C) edge (A);
    \path[onslide={<2>{highlightorange}}] (E) edge (G); 
\end{tikzpicture}
\end{column}
\end{columns}
\end{frame}

\againframe<4>{diameter}

\begin{frame}{Jetzt: Bestimmung des Durchmessers eines Baumes}
\begin{columns}[T]
\begin{column}{0.6\textwidth}
In einen Baum $T = (V, E)$ können wir die \alert{Kreisfreiheit} ausnutzen, denn diese impliziert, dass
\begin{enumerate}
    \item je zwei Knoten stets durch einen \emph{eindeutigen} Pfad miteinander verbunden sind, und
    \item $\abs{E} = \abs{V} - 1$.
\end{enumerate}

\medskip

Den Durchmesser kann man nun mittels zweier \alert{Breitensuchen} bestimmen --- deren Laufzeit liegt f\"ur B\"aume n\"amlich in $\bigO(n)$ f\"ur $n \coloneqq \abs{V}$.

\medskip

Im Folgenden stellen wir zun\"achst den Algorithmus im Detail vor und beweisen anschlie{\ss}end seine Korrektheit sowie die Laufzeitschranke von $\bigO(n)$.
\end{column}
\begin{column}{0.35\textwidth}
\vspace{0.8cm}
\begin{tikzpicture}[every edge/.append style={draw, -}]
    \node[] (A) at (0,0) {A};
    \node[] (B) at (0.5,2) {B};
    \node[] (C) at (2,2.5) {C};
    \node[] (D) at (1.7,-0.2) {D};
    \node[] (E) at (2,-2) {E};
    \node[] (F) at (3,1.3) {F};
    \node[] (G) at (3.2,-1.2) {G};
    
    \path[] (A) edge (B);
    \path[] (A) edge (D);
    \path[] (D) edge (E);
    \path[] (C) edge (F);
    \path[] (C) edge (A);
    \path[] (E) edge (G); 
\end{tikzpicture}
\end{column}
\end{columns}
\end{frame}

\begin{frame}{Algorithmus zur Bestimmung des Durchmessers eines Baumes}
\begin{columns}[T]
\begin{column}{0.5\textwidth}
\vspace{-6pt}
\begin{algorithm}[H]
	\caption{Durchmesser eines Baumes}
	\label{trees:alg:diameter}
	\SetKwFunction{BFS}{BFS}
	\DontPrintSemicolon
	\Input{Baum $T = (V, E)$}
    \Output{Durchmesser von $T$}
    W\"ahle einen beliebigen Knoten $s \in V$\;
	Starte \BFS{$T$, $s$}\;
	Sei $t$ der \emph{zuletzt} entdeckte Knoten\;
	\BlankLine
    Starte \BFS{$T$, $t$}\;
	Sei $u$ der \emph{zuletzt} entdeckte Knoten\;
	\BlankLine
    \Return{L\"ange des Pfades von $t$ nach $u$}
\end{algorithm}
\end{column}
\begin{column}{0.47\textwidth}
\vspace{-6pt}
\begin{remark}<2->
Der durch eine Breitensuche \emph{zuletzt} entdeckte Knoten muss zu den am \emph{weitesten} vom Ausgangsknoten entfernten Knoten geh\"oren.
\end{remark}

\begin{remark}<3->
Aufgrund der Kreisfreiheit von $T$ muss \emph{der} Pfad von $t$ nach $u$ eindeutig sein.
\end{remark}

\begin{theorem}<4->
\label{trees:thm:diameter}
\vspace*{-7pt}
\algref{trees:alg:diameter} ermittelt den Durchmesser eines Baumes in Zeit $\bigO(n)$ f\"ur $n = \abs{V}$.
\end{theorem}
\end{column}
\end{columns}
\end{frame}

\begin{frame}{Beweis von \thmref{trees:thm:diameter}}
\begin{block}<+->{Laufzeit}
Die Laufzeit von Algorithmus \ref{trees:alg:diameter} ist klar, da zum einen f\"ur alle B\"aume $T = (V, E)$ mit $n \coloneqq \abs{V}$ und $m \coloneqq \abs{E}$ gilt, dass $m = n - 1$ ist.
Zudem wissen wir, dass Breitensuche f\"ur alle zusammenh\"angenden Graphen in $\bigO(m)$ liegt --- f\"ur B\"aume also sogar in $\bigO(n)$.
\end{block}

\begin{block}<+->{Korrektheit}
F\"ur den Beweis der Korrektheit sei $T = (V, E)$ ein Baum und bezeichne $\delta(v, w)$ die L\"ange des (eindeutigen) Pfades von $v$ nach $w$ f\"ur alle $v, w \in V$.

\begin{lemma}<+->
\label{trees:lem:dist}
\vspace*{-7pt}
Sei $P = a, \dots, b$ ein l\"angster Pfad in $T$ ($\delta(a, b)$ entspricht also gerade dem Durchmesser von $T$) und o.B.d.A. der \emph{einzige} solche Pfad.
Weiterhin sei $s \in V$ beliebig gew\"ahlt.

Bei einer von $s$ ausgehenden Breitensuche ist entweder $a$ oder $b$ der zuletzt entdeckte Knoten.
\end{lemma}
\end{block}
\end{frame}

\begin{frame}{Beweis von \lemref{trees:lem:dist} \, I}
Wir beweisen per Widerspruch.
Wir nehmen also an, es existierte ein Knoten $x \in V \setminus \{ a, b \}$, der weiter von $s$ entfernt ist, als $a$ und $b$ es sind.
Formal ausgedr\"uckt bedeutete das
\begin{align}\label{trees:eq:ineq}
    \delta(s, a) < \delta(s, x) && \text{und} && \delta(s, b) < \delta(s, x).
\end{align}

Sei nun $c$ derjenige Knoten auf dem Pfad $P = a, \dots, b$, der durch die von $s$ ausgehende Breitensuche zuerst entdeckt wird.
Das ist zugleich der Knoten auf $p_1$, der $s$ am n\"achsten ist.

Durch Einsetzen der ersten Ungleichung aus \eqref{trees:eq:ineq} erhalten wir
\begin{align}\label{trees:eq:f1}
    \underbrace{\delta(s, a)}_{< \delta(s, x)} + \delta(s, b) < \delta(s, x) + \delta(s, b).
\end{align}
\end{frame}

\begin{frame}{Beweis von \lemref{trees:lem:dist} \, II}
Da die Pfade von $s$ nach $a$ und $s$ nach $b$ beide \"uber $c$ f\"uhren, haben sie den Teilpfad von $s$ bis $c$ gemeinsam, bevor sie in unterschiedliche "Richtungen" (gen $a$ bzw. $b$) abzweigen:
\begin{align}\label{trees:eq:eq}
    \delta(s, a) = \delta(s, c) + \delta(c, a) && \text{und} && \delta(s, b) = \delta(s, c) + \delta(c, b).
\end{align}
W\"urden die Gleichungen in \eqref{trees:eq:eq} nicht gelten, m\"usste ein alternativer Pfad von $s$ nach $a$ bzw. $b$ existieren, was die Existenz eines, in B\"aumen per Definition ausgeschlossenen, Kreises bedeuten w\"urde.

Wir erhalten also zus\"atzlich zu der Ungleichung in \eqref{trees:eq:f1} auch folgende Gleichung:
\begin{align}
\begin{split}\label{trees:eq:f2}
    \delta(s, a) + \delta(s, b) &\overset{\eqref{trees:eq:eq}}{=} \delta(s, c) + \delta(c, a) + \delta(s, c) + \delta(c, b)\\
    &= \delta(a, b) + 2 \cdot \delta(s, c).
\end{split}
\end{align}
Zusammen ergeben die Formeln \eqref{trees:eq:f1} und \eqref{trees:eq:f2}
\begin{align}\label{trees:eq:f3}
    \delta(a, b) + 2 \cdot \delta(s, c) < \delta(s, x) + \delta(s, b).
\end{align}
\end{frame}

\begin{frame}{Beweis von \lemref{trees:lem:dist} \, III}
Weiterhin muss gem\"a{\ss} der \emph{Dreiecksungleichung} gelten, dass
\begin{align}\label{trees:eq:f4}
    \delta(x, b) = \delta(s, x) + \delta(s, b).
\end{align}
Andernfalls m\"usste $\delta(x, b) < \delta(s, x) + \delta(s, b)$ sein, was aufgrund der Kreisfreiheit wiederum ausgeschlossen ist.

Aus den Formeln \eqref{trees:eq:f3} und \eqref{trees:eq:f4} und der offensichtlichen Tatsache, dass $\delta(a, b) < \delta(a, b) + 2 \cdot \delta(s, c)$, erhalten wir
\begin{align}
    \delta(a, b) < \delta(a, b) + 2 \cdot \delta(s, c) < \delta(x, b)
\end{align}

Per Voraussetzung sollte aber $P = a, \dots, b$ der \emph{l\"angste} Pfad sein, und nicht $P' = x, \dots, b$ --- Widerspruch! \qed
\end{frame}

\begin{frame}{Abschluss des Beweises von \thmref{trees:thm:diameter}}
Aus \lemref{trees:lem:dist} wissen wir bereits, dass wir durch eine Breitensuche, ausgehend von einem beliebigen Knoten $s$, zuletzt einen Endpunkt $t \in \{ a, b \}$ eines l\"angsten Pfades $P = a, \dots, b$ in $T$ finden.

O.B.d.A. sei $t = a$.
Dann ist klar, dass f\"ur den, bei einer von $t$ ausgehenden Breitensuche zuletzt gefundenen Knoten $u$ gelten muss, dass $u = b$, da dieser am weitesten von $a$ entfernt ist.

Somit entspricht die L\"ange des durch \algref{trees:alg:diameter} gefundenen Pfades tats\"achlich dem Durchmesser von $T$.

Dies schlie{\ss}t den Beweis der Korrektheit von \algref{trees:alg:diameter} und somit auch von \thmref{trees:thm:diameter} ab. \qed
\end{frame}

\begin{frame}{Weitere Eigenschaften von B\"aumen}
\begin{columns}[T]
\begin{column}{0.6\textwidth}
\begin{definition}
\only<1->{Sei $T = (\red<1>{V}, \orange<1>{E}, \purple<1>{r})$ ein gewurzelter Baum.}

\only<2->{Die \alert{Tiefe} von $T$ ist $\orange<2>{d}(T) \coloneqq \max_{\teal<2>{v} \in V}\orange<2>{d(\teal<2>{v})}$.}

\only<3->{Sei $\teal<3>{N^+}(\red<3>{v}) \coloneqq \{ \teal<3>{w} \in V : \orange<3>{\{ v, w \}} \in E, d(\teal<3>{w}) = d(\red<3>{v}) + 1 \}$ die Menge aller \emphteal<3>{Kinder} eines Knotens $\red<3>{v} \in V$.
Dann ergibt sich die \alert{Ordnung} von $T$ zu $\Delta^+ \coloneqq \max_{\red<3>{v} \in V} \abs{\teal<3>{N^+}(\red<3>{v})}$.}

\only<4->{$T$ ist \alert{vollst\"andig}, falls alle Knoten $v \in V$ mit $d(v) < d$ \emph{genau $\Delta^+$} Kinder haben.}

\only<5->{Gewurzelte B\"aume der Ordnung 2 nennen wir \alert{Bin\"arb\"aume}.}
\end{definition}
\end{column}
\begin{column}{0.35\textwidth}
\vspace{0.8cm}
\begin{tikzpicture}[every edge/.append style={draw, -}]
    \tikzstyle{highlightred}=[madderlake, fill=madderlake!7]
    \tikzstyle{highlightteal}=[tealblue, fill=tealblue!7]
    \tikzstyle{highlightorange}=[mandarin]
    \tikzstyle{highlightpurple}=[amaranth, fill=amaranth!7]
    \tikzstyle{transparentblue}=[hublue!30, fill=none]

    \node[onslide={<1, 2>{highlightpurple}}, onslide={<3>{highlightred}}] (A) at (0,0) {A};
    \node[onslide={<1>{highlightred}}, onslide={<3>{highlightteal}}, onslide={<5>{transparentblue}}] (B) at (0.5,2) {B};
    \node[onslide={<1>{highlightred}}, onslide={<3>{highlightteal}}] (C) at (2,2.5) {C};
    \node[onslide={<1, 2>{highlightred}}, onslide={<3>{highlightteal}}] (D) at (1.7,-0.2) {D};
    \node[onslide={<1, 2>{highlightred}}] (E) at (2,-2) {E};
    \node[onslide={<1>{highlightred}}] (F) at (3,1.3) {F};
    \node[onslide={<1>{highlightred}}, onslide={<2>{highlightteal}}] (G) at (3.2,-1.2) {G};
    
    \path[onslide={<1, 3>{highlightorange}}, onslide={<5>{transparentblue}}] (A) edge (B);
    \path[onslide={<1, 2, 3>{highlightorange}}] (A) edge (D);
    \path[onslide={<1, 2>{highlightorange}}] (D) edge (E);
    \path[onslide={<1>{highlightorange}}] (C) edge (F);
    \path[onslide={<1, 3>{highlightorange}}] (C) edge (A);
    \path[onslide={<1, 2>{highlightorange}}] (E) edge (G); 
\end{tikzpicture}
\end{column}
\end{columns}
\end{frame}

\begin{frame}{Suchb\"aume}
\begin{definition}
Ein gewurzelter Baum $T = (V, E, r)$ hei{\ss}t \alert{geordnet}, falls f\"ur alle $v \in V$ die Menge $N^+(v)$ strikt geordnet ist.
\end{definition}

\begin{block}{Vereinbarungen}
\begin{itemize}
    \item Im Folgenden meint "Bin\"arbaum" stets einen \emph{geordneten} Bin\"arbaum.
    \item F\"ur einen inneren Knoten $v$ bezeichnen $v_{\ell}$ und $v_r$ sein "linkes" bzw. "rechtes" Kind, d.h. $v_{\ell} < v_r$.
\end{itemize}
\end{block}

\begin{definition}[Bin\"arer Suchbaum]
Sei $T = (V, E, r)$ ein Bin\"arbaum mit injektiver \alert{Beschriftungsfunktion} $\lambda : V \to \nat$.
Wir nennen $T$ einen \alert{bin\"aren Suchbaum}, falls f\"ur alle Knoten $v \in V$ mit Kindern $v_{\ell}$ und $v_r$ gilt, dass $\lambda(v_{\ell}) < \lambda(v) < \lambda(v_r)$.
\end{definition}
\end{frame}

\begin{frame}{AVL-B\"aume}
\begin{definition}<+->
Sei $T = (V, E, r)$ ein bin\"arer Suchbaum.
F\"ur Knoten $v \in V$ mit Kindern $v_{\ell}, v_r$ seien $d_{\ell} \coloneqq d(T_{v_{\ell}})$ und $d_r \coloneqq d(T_{v_r})$ die H\"ohen des linken und rechten Teilbaumes von $v$.

Die \alert{Balance} von $v$ definieren wir durch $bal(v) \coloneqq \abs{d_{\ell} - d_r}$
\end{definition}

\begin{definition}<+->[AVL-Baum]
Ein AVL-Baum $T = (V, E, r)$ ist ein bin\"arer Suchbaum, in dem für alle $v \in V$ gilt, dass $bal(v) \leq 1$, d.h. die H\"ohen des linken und rechten Teilbaumes von $v$ unterscheiden sich um h\"ochstens eins.

Diese Eigenschaft nennt man \alert{höhenbalanciert}.
\end{definition}

\begin{alertblock}{Achtung}<+->
Das bedeutet \emph{nicht}, dass sich die Tiefen aller Bl\"atter um h\"ochsten eins unterscheiden.
\end{alertblock}
\end{frame}

\begin{frame}{Eigenschaften von AVL-B\"aumen}
\begin{theorem}
F\"ur alle AVL-B\"aume $T = (V, E, r)$ mit $n \coloneqq \abs{V}$ Knoten ist $d(T) \in \bigO(\log n)$.
\end{theorem}
\begin{proof}
Siehe Vorlesung!
\end{proof}
\end{frame}

\begin{frame}{Operationen auf AVL-B\"aumen}
\only<+->{
    Wie f\"ur Suchb\"aume \"ublich, wollen wir mindestens die folgenden Operationen unterst\"utzen:
    \begin{itemize}
        \item Suchen
        \item Einf\"ugen
        \item L\"oschen
    \end{itemize}
}

\only<+->{Im Unterschied zu nat\"urlichen Suchb\"aumen m\"ussen wir hier darauf achten, die H\"ohenbalance zu wahren.}
\begin{itemize}
    \item<+-> Da eine \emph{Suche} nichts an der Form des Baumes \"andert, muss diese auch nicht angepasst werden.
    \item<+-> Anders sieht es beim \emph{Einf\"ugen} und \emph{L\"oschen} aus: Hier gehen wir zun\"achst wie gewohnt vor und stellen die Balance anschlie{\ss}end durch \alert{Rotationen} wieder her.
\end{itemize}
\end{frame}

\begin{frame}{Schreibtischtests zu AVL-B\"aumen}
F\"ur Schreibtischtests zu den eben genannten Operationen auf AVL-B\"aumen sei auf die \alert{\href{https://www2.informatik.hu-berlin.de/~mjung/Tutorium/AlgoDat2016/folien10.pdf}{Folien von Michael R. Jung}} aus dem Jahr 2016 verwiesen.
\end{frame}

    
    \begin{frame}[plain, c]
    \centering\LARGE{ENDE}
    \end{frame}
\end{document}
